\chapter{Infectious Disease}

\medterm{Acute HIV Syndrome}-- The acute retroviral syndrome occurring 2-6 weeks after primary HIV infection, corresponding to high-level viremia and immune activation. Presents with fever (most common), lymphadenopathy, pharyngitis, rash (maculopapular, trunk/face), myalgia, arthralgia, headache, nausea, diarrhea, and oral ulcers. High viral load, low CD4 count transiently. Highly infectious period. Diagnosis: HIV RNA (viral load) is positive; HIV antibody may be negative (window period). p24 antigen is detectable. Treatment: immediate ART regardless of CD4 count; pre-exposure prophylaxis for partners. Early diagnosis and treatment improves outcomes and reduces transmission.

\medterm{AIDS}-- Acquired immunodeficiency syndrome: the advanced stage of HIV infection, defined by CD4 count <200 cells/µL or the presence of AIDS-defining illnesses. Opportunistic infections: Pneumocystis jirovecii pneumonia, candidiasis, toxoplasmosis, cryptococcal meningitis, CMV retinitis, tuberculosis, mycobacterium avium complex; malignancies: Kaposi sarcoma, non-Hodgkin lymphoma, cervical cancer; wasting, dementia. Management: antiretroviral therapy (immune reconstitution), prophylaxis for opportunistic infections (trimethoprim-sulfamethoxazole for PCP when CD4 <200), and treatment of active infections; monitor viral load/CD4. Prevention of progression with early ART dramatically improves prognosis.

\medterm{Amoebiasis}-- Infection caused by the protozoan parasite Entamoeba histolytica, transmitted via the fecal-oral route (contaminated food or water). Most infections are asymptomatic; symptomatic disease ranges from amoebic dysentery (bloody diarrhea, abdominal pain, tenesmus) to extra-intestinal disease, most commonly amoebic liver abscess (right upper quadrant pain, fever, weight loss). Diagnosis: microscopy of stool samples (identification of trophozoites containing ingested red blood cells), antigen detection, PCR, and serology for extra-intestinal disease. Treatment: tinidazole or metronidazole for invasive disease, followed by a luminal agent (paromomycin or diloxanide furoate) to eradicate cyst carriage. Prevention: safe water, hand hygiene, sanitation.

\medterm{Anthrax}-- A serious infectious disease caused by the gram-positive bacterium Bacillus anthracis, which forms highly resistant spores that can persist in the environment for decades, and it is primarily a zoonotic disease of herbivores (cattle, sheep, goats) that can be transmitted to humans through contact with infected animals or their products. In humans, anthrax occurs in three main forms depending on the route of exposure: cutaneous (most common, 95 percent, presenting as a painless pruritic papule that progresses to a vesicle and then to a characteristic black eschar, with surrounding edema), inhalational (the most lethal form, causing hemorrhagic mediastinitis and mediastinal widening on chest radiography, with a biphasic illness of initial nonspecific symptoms followed by sudden respiratory deterioration), and gastrointestinal (following ingestion of contaminated meat, causing abdominal pain, nausea, vomiting, and bloody diarrhea). Diagnosis is confirmed by culture, Gram stain, or PCR from blood, skin lesions, or tissue samples, and serology may be used retrospectively. Treatment involves high-dose IV antibiotics: ciprofloxacin or doxycycline in combination with one or two additional agents (clindamycin, rifampin, meropenem), and antitoxin (raxibacumab or obiltoxaximab) may be considered for inhalational anthrax. Post-exposure prophylaxis consists of ciprofloxacin or doxycycline for 60 days combined with anthrax vaccine adsorbed (AVA, BioThrax).

\medterm{Antibiotic}-- A substance that kills or slows the growth of bacteria.

\medterm{Antimicrobial}-- A general term for antibiotics and other drugs that fight microscopic organisms in the body, such as bacteria, viruses, fungi, and parasites.

\medterm{Antimicrobial Resistance Mechanisms}-- Beta-lactamases: Ambler class A (TEM, SHV, CTX-M ESBLs, KPC carbapenemases), class B (metallo-beta-lactamases: NDM, VIM, IMP zinc-dependent, not inhibited by avibactam), class C (AmpC, chromosomal/inducible, hydrolyzes cephalosporins), class D (OXA carbapenemases,.

\medterm{Antimicrobial Stewardship}-- Core elements (CDC): leadership commitment, accountability, drug expertise, action (prospective audit/feedback, preauthorization), tracking (DOT,.

\medterm{Antiseptic}-- Substances used on wounds to prevent or treat infection; they kill or slow the growth of diseasecausing organisms, such as bacteria, on the surface of the body.

\medterm{Antiviral}-- A class of drugs that inhibit viral replication, used to treat viral infections. Mechanisms: nucleoside/nucleotide analogs (acyclovir—HSV/VZV, ganciclovir—CMV, tenofovir/lamivudine—HIV/HBV, remdesivir—RNA viruses), protease inhibitors (HIV), neuraminidase inhibitors (oseltamivir, zanamivir—influenza), polymerase inhibitors, integrase inhibitors, entry/fusion inhibitors, and monoclonal antibodies. Used for herpesviruses, HIV (ART), hepatitis B/C, influenza, COVID-19, RSV (palivizumab prophylaxis), and CMV. Adverse effects: nephrotoxicity, myelosuppression, GI effects, and resistance with incomplete suppression. Aimed at reducing viral load, symptoms, transmission, and complications.

\medterm{Arbovirus}-- A virus transmitted by mosquitoes or other member of the arthropod phylum.

\medterm{Ascariasis}-- Infection with Ascaris lumbricoides, the most common helminth infection worldwide (soil-transmitted). Lifecycle: ingestion of embryonated eggs $\rightarrow$ larvae hatch in small intestine $\rightarrow$ penetrate intestinal wall $\rightarrow$ migrate through liver and lungs (Loeffler syndrome: cough, eosinophilia, transient pulmonary infiltrates) $\rightarrow$ return to small intestine as adults. Most infections are asymptomatic. Heavy burden: intestinal obstruction (especially children), biliary obstruction, pancreatitis, malnutrition in children. Diagnosis: stool microscopy (characteristic eggs). Treatment: albendazole or mebendazole (single dose).

\medterm{Aspergillosis}-- A spectrum of diseases caused by the fungus Aspergillus (most commonly A. fumigatus). Allergic forms: allergic bronchopulmonary aspergillosis (ABPA, in asthma or cystic fibrosis -- wheezing, eosinophilia, pulmonary infiltrates), allergic Aspergillus sinusitis. Saprophytic: aspergilloma (fungus ball in pre-existing lung cavity -- haemoptysis). Invasive aspergillosis: occurs in immunocompromised patients (neutropenia, transplant, steroid therapy) -- hyphal invasion of blood vessels causes tissue infarction, pneumonia, and dissemination to brain, heart, and kidneys (high mortality). Diagnosis: serum galactomannan, beta-D-glucan, CT chest (halo sign, air crescent sign), culture, biopsy. Treatment: voriconazole (first-line), amphotericin B, isavuconazole.

\medterm{Athlete Foot}-- Tinea pedis: dermatophyte infection (Trichophyton rubrum, T. mentagrophytes) of the feet, common in those who wear occlusive footwear or use communal showers. Forms: interdigital (scaly, fissured, macerated between toes—most common), moccasin (diffuse plantar scaling), vesicular. Features: itching, burning, cracking, odor; predisposes to bacterial cellulitis (portal of entry). Diagnosis: clinical; KOH preparation, culture. Management: topical antifungals (terbinafine, clotrimazole) for 2-4 weeks; oral terbinafine/itraconazole for resistant or extensive cases; hygiene (dry feet, clean socks, footwear rotation, avoid walking barefoot in public).

\medterm{Bacillary Angiomatosis}-- A vascular proliferative disease caused by Bartonella henselae or B. quintana, occurring in immunocompromised patients (HIV/AIDS, transplant). Presents with erythematous papules, nodules, or pedunculated lesions resembling Kaposi sarcoma or pyogenic granuloma. Can involve skin, liver (peliosis hepatis), spleen, bone, and lymph nodes. Fever, night sweats, and weight loss may accompany. Diagnosis: biopsy shows lobular proliferations of blood vessels with epithelioid endothelial cells and Warthin-Starry silver stain positive bacilli; PCR confirmation. Treatment: erythromycin or doxycycline for 3 months. Lesions resolve with treatment.

\medterm{Bacteremia}-- The presence of viable bacteria in the bloodstream, a serious condition that can be transient (after dental procedures, toothbrushing), intermittent (undrained abscess), or continuous (endovascular infection like endocarditis). Bacteremia triggers systemic inflammatory responses ranging from fever and leukocytosis to septic shock. Diagnosis requires positive blood cultures; treatment involves IV antibiotics targeting the identified organism and source control of the primary infection site.

\medterm{Bacterial Infections}-- Illnesses caused by pathogenic bacteria proliferating in the body, ranging from mild sore throats and ear infections to life-threatening sepsis and meningitis. They are diagnosed through cultures, PCR testing, or other microbiological methods. Treatment typically involves antibiotics, though rising antimicrobial resistance demands careful stewardship.

\medterm{Bartonellosis}-- A group of infections caused by Bartonella species. Cat scratch disease (B. henselae): regional lymphadenopathy after cat scratch/bite, fever, inoculation papule. Bacillary angiomatosis (B. henselae, B. quintana): vascular proliferative lesions in immunocompromised (HIV), red papules/nodules on skin, bone, and viscera. Oroya fever (B. bacilliformis, endemic in Peru): acute hemolytic anemia, fever, severe myalgia (Carrion disease). Verruga peruana: chronic cutaneous nodular stage. Trench fever (B. quintana): relapsing fever, bone pain, louse-borne. Diagnosis: culture, PCR, serology. Treatment: azithromycin for CSD, doxycycline + rifampin for bacillary angiomatosis.

\medterm{Bioterrorism Agents}-- Category A: anthrax (Bacillus anthracis), smallpox (variola), botulism (Clostridium botulinum), plague (Yersinia pestis), tularemia (Francisella tularensis), viral hemorrhagic fevers (Ebola, Marburg, Lassa, Machupo). Anthrax: inhalational (mediastinal widening, hemorrhagic mediastinitis), cutaneous (eschar), GI. Treatment: ciprofloxacin/doxycycline x60d + vaccine (BioThrax). Smallpox: eradicated 1980, vaccine (ACAM2000, Jynneos), tecovirimat, brincidofovir. Botulism: descending flaccid paralysis, antitoxin (heptavalent). Plague: bubonic (lymphadenopathy/bubo), pneumonic (person-to-person), septicemic. Doxycycline/ciprofloxacin/gentamicin. Tularemia: ulceroglandular, oculoglandular, pneumonic, typhoidal. Streptomycin/gentamicin/doxycycline. VHF: supportive, specific antivirals where available. Public health response: isolation, contact tracing, prophylaxis, risk communication.

\medterm{Blastomycosis}-- A fungal infection caused by Blastomyces dermatitidis, endemic to the Ohio and Mississippi River valleys and Great Lakes region. Inhalation of conidia causes pulmonary infection. Pulmonary: acute pneumonia (often misdiagnosed as bacterial), chronic pneumonia with nodules and cavitation. Disseminated: skin (verrucous or ulcerative lesions, most common extrapulmonary site), bone (osteomyelitis), and GU (prostatitis, epididymitis). Diagnosis: sputum culture, histopathology (broad-based budding yeasts), serum antigen testing. Treatment: itraconazole for mild-moderate; amphotericin B for severe or CNS disease.

\medterm{Borrelia}-- A genus of spiral-shaped bacteria (spirochetes) transmitted by arthropod vectors, most notably ticks and lice. The major human pathogen is Borrelia burgdorferi sensu lato, the causative agent of Lyme disease, transmitted by Ixodes ticks. Borrelia recurrentis causes louse-borne relapsing fever, and Borrelia duttonii causes tick-borne relapsing fever. Borrelia species have a complex genome with multiple linear and circular plasmids. Diagnosis relies on serology (ELISA with Western blot confirmation), PCR, or direct microscopy. Treatment: doxycycline, amoxicillin, or ceftriaxone depending on disease stage.

\medterm{Botulism}-- A rare, potentially fatal neuroparalytic disease caused by botulinum toxin, a potent presynaptic neurotoxin produced by Clostridium botulinum, an anaerobic, spore-forming, gram-positive bacillus. Botulinum toxin blocks the release of acetylcholine at the neuromuscular junction by cleaving SNARE proteins (SNAP-25, VAMP/synaptobrevin, syntaxin), preventing vesicle fusion and causing flaccid paralysis. Four clinical forms exist: foodborne botulism (ingestion of preformed toxin in improperly canned or preserved foods, most commonly homecanned vegetables, honey, and fermented foods), infant botulism (the most common form in the United States, caused by intestinal colonization and in vivo toxin production after ingestion of C. botulinum spores, most commonly from honey), wound bo-.

\medterm{Brucellosis}-- A zoonotic infectious disease caused by gram-negative coccobacilli of the genus Brucella, transmitted to humans through contact with infected animals (cattle, goats, pigs, sheep) or consumption of unpasteurized dairy products. Brucella melitensis (goats and sheep, most severe), B. abortus (cattle), B. suis (pigs), and B. canis (dogs) are the principal species. It is endemic in the Mediterranean basin, Middle East, Central Asia, Mexico, and South America. Clinical presentation is highly variable: acute brucellosis presents with undulant fever (fever that rises and falls in waves), sweats (often malodorous), arthralgia, fatigue, headache, and hepatosplenomegaly. Chronic brucellosis may persist for months to years with relapsing fever, chronic fatigue, depression, arthritic manifestations.

\medterm{Candidiasis}-- Fungal infection caused by Candida species, most commonly C. albicans. Mucocutaneous: oral thrush (white plaques on oral mucosa), esophagitis, vulvovaginitis (itching, discharge), and intertrigo. Invasive candidiasis: candidemia and deep organ infection, typically in immunocompromised or ICU patients. Diagnosis: culture, histopathology. Treatment: topical or oral azoles for mucocutaneous; echinocandins for invasive disease. Risk factors: antibiotics, immunosuppression, diabetes, indwelling catheters.

\medterm{Cardiac Device Infections}-- CIED (pacemaker, ICD) infection: pocket infection (erythema, warmth, purulence, erosion), lead infection (bacteremia, endocarditis). Diagnosis: blood cultures, TEE (lead vegetations), pocket culture. Treatment: complete system removal (leads + generator) - CLASS I. Extraction: transvenous (laser, mechanical sheaths) vs surgical (thoracotomy). Antibiotics: pocket only - 10-14d; bacteremia/endocarditis 4-6w. Relapsing bacteremia if hardware retained. Bridge therapy: temporary transvenous pacing, wearable defibrillator (LifeVest). Reimplantation: contralateral side after clearance.

\medterm{Cat Scratch Disease}-- A regional lymphadenitis caused by Bartonella henselae, a gramnegative bacillus transmitted to humans through cat scratches, bites, or licks on broken skin, particularly from kittens (which carry higher bacterial loads). Clinical presentation: 1-2 weeks after exposure, a papule or pustule develops at the inoculation site, followed by regional lymphadenopathy (axillary, epitrochlear, cervical) that is tender and may suppurate (10-15 percent). Systemic symptoms include low-grade fever, fatigue, headache, anorexia. Atypical presentations: Parinaud oculoglandular syndrome (conjunctivitis + preauricular lymphadenopathy), encephalitis, neuroretinitis, osteomyelitis, hepatosplenic granulomas, endocarditis (culture-negative). Diagnosis: Warthin-Starry silver.

\medterm{Cat-Scratch Fever}-- An infectious disease caused by the bacterium Bartonella henselae, transmitted through a scratch, bite, or lick from an infected cat (typically kittens). After an incubation period of 3--14 days, a papule or pustule forms at the inoculation site, followed by painful regional lymphadenopathy (axillary, epitrochlear, cervical, or submandibular) within 1--3 weeks. Most cases are self-limited, resolving over 2--4 months. Atypical presentations: Parinaud oculoglandular syndrome (conjunctivitis with preauricular lymphadenopathy), hepatosplenic disease (fever, abdominal pain), neuroretinitis, and encephalopathy. Diagnosis: serology (IgG by IFA or ELISA), PCR of lymph node aspirate, or histology (Warthin-Starry silver stain shows pleomorphic bacilli). Treatment: typically supportive; azithromycin for moderate-to-severe lymphadenopathy or immunocompromised patients.

\medterm{Cell-mediated immunity}-- A type of immune response mounted against viruses, certain types of parasites, and perhaps cancer cells.

\medterm{Cellulitis}-- A bacterial infection of the deep dermis and subcutaneous tissue, most commonly caused by Streptococcus pyogenes (Group A Strep) and Staphylococcus aureus. Presents with erythema, warmth, edema, and tenderness, with poorly defined margins. Often on the lower extremities. Risk factors: breaks in skin (trauma, ulcers, tinea pedis), lymphedema, venous insufficiency, and immunosuppression. Diagnosis: clinical; blood cultures if systemic signs. Treatment: antibiotics targeting Strep and Staph (cephalexin, dicloxacillin, clindamycin). Severe cases or facial involvement require IV antibiotics.

\medterm{Chancroid}-- A sexually transmitted infection caused by Haemophilus ducreyi, characterized by painful genital ulcers and tender inguinal lymphadenopathy (bubo). Ulcers are ragged, non-indurated, with purulent exudate and surrounding erythema. More common in men and in resource-limited settings. Associated with increased HIV transmission. Diagnosis: clinical (painful ulcer + tender lymphadenopathy), Gram stain (school of fish appearance), culture (chocolate agar with vancomycin). Treatment: azithromycin (single dose), ceftriaxone (single dose IM), ciprofloxacin, or erythromycin. Fluctuant buboes require aspiration.

\medterm{Cholera}-- An acute diarrheal illness caused by Vibrio cholerae, a gram-negative, comma-shaped, flagellated bacterium, serogroups O1 and O139, that produces cholera toxin, an enterotoxin that irreversibly activates adenylate cyclase in intestinal epithelial cells, causing massive secretion of isotonic fluid into the intestinal lumen. Transmission is fecal-oral through contaminated water and food, with humans as the only natural host. Seven pandemics have occurred since 1817, and cholera remains endemic in over 50 countries, particularly in regions with inadequate water and sanitation. The hallmark presentation is acute onset of profuse, painless, watery diarrhea described as rice-water stool (cloudy white fluid with flecks of mucus) that can exceed one liter per hour, leading to severe dehydration, metabolic acidosis, and hypovolemic shock within hours of onset. Vomiting may precede diarrhea. Clinical dehydration is classified as severe (two or more signs: delayed capillary refill, weak or absent radial pulse, sunken eyes, decreased skin turgor, altered consciousness) or some dehydration (less severe signs). Diagnosis is clinical; stool culture on thiosulfate-citrate-bile salts-sucrose (TCBS) agar grows yellow colonies, and rapid dipstick tests detect V. cholerae O1 and O139 antigen. Treatment is primarily fluid replacement: WHO low-osmolality oral rehydration solution (ORS: 75 mmol/L glucose, 75 mEq/L sodium) for some or no dehydration; aggressive IV Ringer’s lactate for severe dehydration (100 mL/kg over 3 hours for adults, 100 mL/kg over 6 hours for children). Antibiotics reduce the duration and volume of diarrhea: azithromycin 1 gram orally as a single dose (preferred for all ages, including pregnant.

\medterm{Clostridioides difficile Infection}-- See gastroenterology.

\medterm{Clostridium Difficile}-- Clostridioides difficile infection (CDI): a toxin-producing anaerobic gram-positive rod causing antibiotic-associated diarrhea and pseudomembranous colitis, most commonly after broad-spectrum antibiotics (clindamycin, cephalosporins, fluoroquinolones). Features: watery diarrhea (often >3 stools/day), abdominal cramping, fever; severe: ileus, toxic megacolon, perforation. Diagnosis: GDH + toxin EIA or toxigenic PCR; consider when diarrhea in hospitalized/antibiotic-exposed patients. Management: stop culprit antibiotics, metronidazole (mild) or oral vancomycin/fidaxomicin (moderate-severe), and fecal microbiota transplant for recurrent; infection control (handwashing, contact precautions, environmental cleaning).

\medterm{CMV}-- See Cytomegalovirus Infection.

\medterm{Coccidioidomycosis}-- A fungal infection caused by Coccidioides immitis or C. posadasii, endemic to the southwestern US (San Joaquin Valley fever). Inhalation of arthroconidia causes primary pulmonary infection. Most cases are asymptomatic or mild (flu-like illness with cough, fever, erythema nodosum). Disseminated disease occurs in immunocompromised patients, affecting skin, bone/joints, and meninges. Diagnosis: serology (IgM, IgG), culture, and histopathology (spherules). Treatment: fluconazole or itraconazole; amphotericin B for severe disseminated disease.

\medterm{Cold Sore}-- Herpes labialis: recurrent herpes simplex virus type 1 (HSV-1) infection of the lips/around mouth, presenting as prodromal tingling, then grouped vesicles that crust and heal in 7-14 days. Triggered by stress, fever, sunlight, immunosuppression, illness. Recurrence due to reactivation of latent virus in the trigeminal ganglion. Diagnosis: clinical; PCR or culture if atypical. Management: topical antivirals (acyclovir cream) or oral antivirals (valacyclovir, famciclovir) initiated at prodrome reduce duration; sunscreen and trigger avoidance; severe/immunocompromised cases need systemic therapy. Transmission by contact with lesions.

\medterm{Common Cold}-- An acute, self-limited viral infection of the upper respiratory tract, one of the most frequent illnesses in humans, with adults experiencing 2-4 episodes annually and children 6-8 episodes. Over 200 viral serotypes cause the common cold, with rhinoviruses (100 serotypes, most common, accounting for 3050 percent of cases) as the leading cause, followed by coronaviruses (15-30 percent, including seasonal HCoV-OC43, -229E, -NL63, -HKU1), respiratory syncytial virus (RSV, 5-15 percent), parainfluenza virus (5-10 percent), adenovirus (5 percent), and enterovirus. Transmission occurs through direct contact, fomites, and respiratory droplets, with peak infectivity in the first 2-3 days of symptoms. The incubation period is 1-3 days. Pathophysiology: viral entry through the nasal epithelium triggers an innate immune response with release.

\medterm{Communicable disease}-- Any disease caused by bacteria, viruses, or other pathogens that is spread from person to person.

\medterm{Community-Acquired Pneumonia}-- CAP: acquired outside hospital. CURB-65 (confusion, uremia >7, respiratory rate >=30, BP <90/60, age >=65) and PSI/PORT score for severity/site of care. Outpatient: amoxicillin 1g TID, doxycycline 100 mg BID, or macrolide (if resistance <25.

\medterm{COVID-19}-- Coronavirus disease 2019 caused by SARS-CoV-2, first identified in December 2019. Ranges from asymptomatic to severe respiratory failure and death. Common symptoms: fever, cough, dyspnea, loss of taste/smell (anosmia), myalgia, and fatigue. Severe disease: pneumonia, ARDS, hyperinflammatory state, and thromboembolic complications. Diagnosis: RT-PCR, rapid antigen testing. Prevention: vaccination (mRNA, adenoviral vector). Treatment: supportive care, remdesivir, dexamethasone for severe disease, and anticoagulation.

\medterm{Cryptococcosis}-- A fungal infection caused by Cryptococcus neoformans or C. gattii, primarily affecting immunocompromised patients (HIV/AIDS, organ transplant, hematologic malignancy). Pulmonary cryptococcosis: cough, dyspnea, pulmonary nodules. Meningoencephalitis: fever, headache, altered mental status, and elevated intracranial pressure. Diagnosis: CSF India ink stain, cryptococcal antigen (CrAg), culture. Treatment: amphotericin B + flucytosine for induction, fluconazole for maintenance. Elevated ICP requires therapeutic lumbar punctures.

\medterm{Cryptosporidiosis}-- A diarrhoeal disease caused by the protozoan parasite Cryptosporidium parvum (zoonotic) or C. hominis (anthroponotic), transmitted via contaminated water (chlorine-resistant oocysts), food, or direct contact. Presents with watery diarrhea, abdominal cramps, nausea, and low-grade fever, lasting 1-2 weeks. In immunocompromised patients (especially HIV with low CD4), it causes chronic, profuse, life-threatening diarrhea and wasting. Diagnosis: modified acid-fast stain of stool (oocysts appear red), antigen detection, PCR. Treatment: supportive (rehydration, electrolyte replacement). Nitazoxanide is effective in immunocompetent patients.

\medterm{Cystitis}-- Inflammation of the urinary bladder, typically from bacterial infection, one of the most common bacterial infections in women. Uncomplicated cystitis occurs in otherwise healthy, non-pregnant women with normal urinary tract anatomy; complicated cystitis occurs in men, pregnant women, children, patients with urinary tract abnormalities, indwelling catheters, immunosuppression, or recent instrumentation. The most common causative organism is Escherichia coli (75-95 percent of uncomplicated cases), followed by Staphylococcus saprophyticus (5-15 percent, particularly in young sexually active women), Klebsiella pneumoniae, and Proteus mirabilis. Clinical presentation includes urinary frequency, urgency, dysuria (painful urination), suprapubic pain, and occasionally hematuria. Fever and flank pain suggest upper tract involvement (pyelonephritis). Diagnosis: urinalysis shows pyuria (leukocyte esterase positive, greater than 5 WBCs per high-power field), nitrite positivity (gram-negative bacteria, not S. saprophyticus or enterococci), and bacteriuria. Urine culture (greater than 103 CFU/mL of a single uropathogen for symptomatic women) confirms the diagnosis and provides antibiotic sensitivity. Treatment: uncomplicated cystitis in women is treated with a short course (3 days) of nitrofurantoin (100 mg twice daily), trimethoprim-sulfamethoxazole (160/800 mg twice daily, if local resistance less than 20 percent), or fosfomycin (3 gram single dose). Beta-lactams (cephalexin, amoxicillin-clavulanate) are less effective and require longer duration (5-7 days). Complicated cystitis requires 7-14 days of culture-directed antibiotics. Men with cystitis require longer treatment (7-14 days) and evaluation for underlying prostatitis or obstruction. Recurrent cystitis (3 or more episodes per year) may be managed with prophylactic antibiotics (postcoital or continuous daily), behavioral modifications (hydration, postcoital voiding), and vaginal estrogen.

\medterm{Cytomegalovirus Infection}-- A widespread herpesvirus infection. Primary infection in immunocompetent individuals is usually asymptomatic or causes mononucleosis-like syndrome. Congenital CMV: most common congenital infection, causes sensorineural hearing loss, microcephaly, and developmental delay. In immunocompromised (HIV/AIDS, transplant): retinitis, colitis, esophagitis, pneumonitis, and encephalitis. Diagnosis: CMV PCR, pp65 antigen, serology. Treatment: ganciclovir, valganciclovir, foscarnet for resistant cases. Prevention: valganciclovir prophylaxis in high-risk transplant recipients.

\medterm{Dengue fever}-- An acute mosquito-borne viral illness caused by dengue virus (DENV-1, -2, -3, -4), a flavivirus transmitted primarily by Aedes aegypti mosquitoes. Dengue is the most rapidly spreading mosquitoborne viral disease worldwide, with an estimated 390 million infections annually and 100 million symptomatic cases, endemic in over 100 countries, particularly in Southeast Asia, the Western Pacific, the Americas, and Africa. The clinical course has three phases: febrile phase (abrupt onset of high fever lasting 2-7 days, severe headache, retro-orbital pain, myalgia, arthralgia, nausea, vomiting, and a characteristic maculopapular rash with islands of sparing); critical phase (occurs around defervescence at days 3-7, lasting 24-48 hours, marked by increased capillary permeability, plasma leakage, and risk of progression to dengue hemorrhagic fever and dengue shock syndrome, characterized by rising hematocrit with falling platelets, serous effusions, and hypoproteinemia); and recovery phase (reabsorption of extravasated fluid, diuresis, and clinical improvement). The World Health Organization classification categorizes severity as dengue without warning signs, dengue with warning signs (abdominal pain, persistent vomiting, fluid accumulation, mucosal bleeding, lethargy, liver enlargement greater than 2 cm, rising hematocrit with rapid platelet drop), and severe dengue (severe plasma leakage leading to shock, severe bleeding, or severe organ involvement). Laboratory findings include leukopenia, thrombocytopenia, elevated liver enzymes, and prolonged aPTT. Diagnosis is confirmed by NS1 antigen detection (positive in the first 5 days), RT-PCR (days 1-5), or serology (IgM positive after days 5-7, IgG rises after 14 days). Treatment is supportive: fever control with acetaminophen (avoid NSAIDs and aspirin due to bleeding risk), aggressive IV fluid resuscitation for severe dengue with plasma leakage (crystalloids, colloids for hypotensive shock), blood transfusion for significant hemorrhage, and close monitoring of vital signs, hematocrit, and platelet counts during the critical phase. No specific an334.

\medterm{Diphtheria}-- An acute, toxin-mediated infectious disease caused by Corynebacterium diphtheriae, a gram-positive aerobic bacillus that produces a potent exotoxin that inhibits protein synthesis via ADP-ribosylation of elongation factor 2, causing local tissue necrosis and systemic toxicity. The bacteria are transmitted through respiratory droplets or direct contact with cutaneous lesions. The classic clinical presentation is respiratory diphtheria: after a 2-5 day incubation period, the patient develops low-grade fever, sore throat, malaise, and a characteristic grayish-white pseudomembrane that forms on the tonsils, pharynx, or larynx composed of fibrin, necrotic epithelial cells, and inflammatory cells. The membrane is firmly adherent, bleeds upon attempted removal, and may extend to obstruct the airway causing stridor, respiratory distress, and asphyxiation. Cervical lymphadenopathy and severe neck edema produce the classic bull neck appearance. Systemic effects of the toxin include myocarditis (occurring in 10-25 percent of cases, presenting with ECG changes, arrhythmias, and heart failure typically 1-2 weeks after onset) and neurologic involvement (palatal paralysis, cranial neuropathies, and peripheral neuritis developing weeks later). Cutaneous diphtheria presents as non-healing ulcers with a gray membrane, more common in tropical settings and individuals with poor hygiene. Diagnosis is clinical confirmed by culture on tellurite or Loeffler medium and the Elek test for toxigenicity. Treatment: equine diphtheria antitoxin (DAT) is the mainstay, neutralizing unbound toxin; it must be given early, intravenously for respiratory cases, with monitoring for serum sickness. Antibiotics eradicate the bacteria and stop toxin production: penicillin G (procaine or IV) or erythromycin for 14 days. Close contacts require throat cultures, prophylaxis with penicillin or erythromycin, and booster immunization. Prevention is through routine childhood immunization with the DTaP vaccine, with boosters (Td or Tdap) every 10.

\medterm{Donovanosis}-- Also known as granuloma inguinale, a chronic, progressive, ulcerative genital infection caused by the intracellular gram-negative bacterium Klebsiella granulomatis (formerly Calymmatobacterium granulomatis). The disease is rare in developed countries but endemic in tropical and subtropical regions, including Papua New Guinea, parts of India, southern Africa, the Caribbean, and South America. Transmission occurs through direct sexual contact with an active ulcer, and the incubation period ranges from 1 week to 3 months after exposure. Clinical presentation begins as a painless, erythematous papule or nodule in the genital region (glans penis, prepuce, labia, vaginal mucosa, cervix, or perianal area) that progressively ulcerates into a painless, beefy-red, friable, granulating ulcer with rolled edges that bleeds easily on contact without significant lymphadenopathy. The ulcers are highly vascular and may have a characteristic beefy red appearance. Without treatment, the ulcer spreads slowly and may cause extensive destruction of genital tissue, fibrosis, lymphedema, and elephantiasis (pseudo-elephantiasis from lymphatic obstruction). Extragenital involvement may occur through autoinoculation to the mouth, lips, and anus. Diagnosis is clinical and confirmed by demonstrating.

\medterm{Dysentery}-- See gastroenterology.

\medterm{Ebola}-- Ebola virus disease (EVD): a severe, often fatal viral hemorrhagic fever caused by Ebola virus (Filoviridae, RNA), transmitted by direct contact with blood/body fluids of infected humans/animals. Incubation 2-21 days. Features: sudden fever, severe headache, myalgia, vomiting, diarrhea, then hemorrhagic manifestations (petechiae, bleeding), multi-organ failure, shock; high case fatality (25-90 percent). Diagnosis: RT-PCR, antigen tests; biocontainment required. Management: supportive (fluids, electrolytes), experimental therapies (monoclonal antibodies—REGEN-EB, remdesivir); infection control (isolation, PPE) is critical; ring vaccination (rVSV-ZEBOV) effective. Endemic to Central/West Africa.

\medterm{Ebola Virus Disease}-- A severe, often fatal hemorrhagic fever caused by the Ebola virus (Filoviridae family). Transmitted through direct contact with blood or body fluids of infected humans or animals. Incubation 2-21 days. Presents with acute onset of fever, myalgia, headache, followed by vomiting, diarrhea, and hemorrhagic manifestations (petechiae, ecchymosis, bleeding). High case fatality rate. Diagnosis: RT-PCR. Treatment: supportive care, monoclonal antibodies. Outbreak control: isolation, contact tracing, safe burials.

\medterm{E Coli Infection}-- Infections caused by Escherichia coli, a gram-negative enteric bacillus. Causes: UTI (most common—cystitis, pyelonephritis), diarrhea (enterotoxigenic ETEC—traveler's diarrhea; enteropathogenic; enterohemorrhagic EHEC O157:H7—bloody diarrhea, HUS; enteroinvasive), neonatal meningitis, sepsis, and intra-abdominal infections. Diagnosis: culture (urine, stool, blood), toxin testing (EHEC). Management: UTI—antibiotics guided by sensitivity; EHEC—avoid antibiotics (increase HUS risk), supportive care; sepsis—empiric broad-spectrum (third-generation cephalosporin, aminoglycoside). Multidrug-resistant (ESBL, carbapenem-resistant) strains complicate therapy; strict hygiene.

\medterm{Ehrlichiosis}-- A tick-borne bacterial infection caused by Ehrlichia chaffeensis (human monocytic ehrlichiosis) and E. ewingii (granulocytic). Endemic in the southeastern and south-central US. Presents with fever, headache, myalgia, malaise, and thrombocytopenia/leukopenia. Rash less common than RMSF. Diagnosis: clinical suspicion in tick-endemic area, PCR of blood, serology (IFA), intracytoplasmic morulae on peripheral smear. Treatment: doxycycline (first-line). Untreated can be severe with respiratory failure, meningoencephalitis, and DIC.

\medterm{Elephantiasis}-- Massive enlargement of a body part, typically a limb or the scrotum, caused by chronic lymphedema from obstruction of lymphatic vessels. The most common cause worldwide is lymphatic filariasis, a mosquito-borne parasitic infection by Wuchereria bancrofti, Brugia malayi, or Brugia timori. Adult filarial worms inhabit the lymphatic vessels, causing inflammation, dilation, and eventual fibrosis and obstruction of lymph flow. Recurrent secondary bacterial infections (cellulitis, lymphangitis) worsen the edema. The affected limb becomes grossly swollen with thickened, hyperkeratotic, verrucous skin. Other causes: non-filarial elephantiasis (podoconiosis) from chronic exposure to irritant volcanic soil in barefoot agricultural workers, post-surgical lymphedema (axillary node dissection, inguinal node dissection), and congenital lymphatic malformations. Treatment: elevation, compression, meticulous skin hygiene, antibiotics for superinfection, and surgical debulking in severe cases.

\medterm{Endemic}-- Continually present among people in a geographic region.

\medterm{Endocarditis Prophylaxis}-- Infective endocarditis prophylaxis with antibiotics is recommended for patients at highest risk of adverse outcomes from endocarditis who are undergoing dental procedures that involve manipulation of gingival tissue, the periapical region of teeth, or perforation of the oral mucosa. The American Heart.

\medterm{Epstein-Barr Virus}-- EBV (human herpesvirus 4): ubiquitous virus causing infectious mononucleosis, transmitted via saliva ('kissing disease'). Features: fever, tonsillopharyngitis, cervical lymphadenopathy, fatigue, splenomegaly; atypical lymphocytosis; complications: splenic rupture (avoid contact sport), hepatitis, hemolytic anemia, airway obstruction; association with malignancies: Burkitt lymphoma, nasopharyngeal carcinoma, Hodgkin lymphoma, post-transplant lymphoproliferative disease; chronic active EBV, hairy leukoplakia in HIV. Diagnosis: heterophile (Monospot) antibody, EBV serology (VCA IgM/IgG, EBNA), PCR. Management: supportive; antivirals not effective; avoid contact sports 4-6 weeks.

\medterm{Febrile}-- Feverish; having a high body temperature.

\medterm{Flu}-- Influenza: a highly contagious acute respiratory illness caused by influenza A and B viruses (orthomyxovirus, RNA, frequent antigenic drift/shift). Features: sudden fever, myalgia, headache, cough, sore throat, malaise; complications: pneumonia (viral or bacterial—S. aureus, S. pneumoniae), exacerbation of COPD/asthma, myocarditis, encephalitis; high risk: elderly, young, pregnancy, chronic disease, immunocompromised. Diagnosis: clinical; rapid antigen/RT-PCR. Management: supportive; antivirals (oseltamivir, zanamivir) within 48 h for high-risk/severe; annual vaccination (inactivated/live attenuated); infection control (cough etiquette, isolation). Pandemic potential via antigenic shift.

\medterm{Fungal Infections}-- Candidiasis: mucosal (thrush, esophagitis, vulvovaginal) - fluconazole; invasive (candidemia, intraabdominal, UTI) - echinocandin (caspofungin, micafungin, anidulafungin) first-line, step-down to fluconazole if susceptible. Aspergillosis: allergic bronchopulmonary (ABPA, asthma/CF, corticosteroids + itraconazole), chronic pulmonary (CPA, cavitary), invasive (IA, neutropenia, transplant, steroids) voriconazole first-line, isavuconazole alternative. Mucormycosis: rhinocerebral (DKA, neutropenia), pulmonary, cutaneous. Amphotericin B lipid formulation (liposomal 5 mg/kg) + surgical debridement; posaconazole/isavuconazole step-down. Cryptococcosis: meningoencephalitis (HIV CD4 <100), pulmonary. Induction: amphotericin B + flucytosine x2w, consolidation: fluconazole x8w, maintenance: fluconazole x1yr. Endemic mycoses: Histoplasma (Ohio/Mississippi valleys), Coccidioides (Southwest US, Valley fever), Blastomyces (Great Lakes, Southeast), Paracoccidioides (Latin America). Treatment: mild-moderate - itraconazole; severe - amphotericin B then itraconazole. Pneumocystis jirovecii pneumonia (PJP): HIV CD4 <200, immunosuppressed. TMP-SMX (first-line), alternatives: dapsone+TMP, atovaquone, clindamycin+primaquine, pentamidine. Prophylaxis: TMP-SMX if CD4 <200.

\medterm{Genital Herpes}-- Herpes simplex virus infection of the genital region (HSV-2 most commonly; HSV-1 increasingly), transmitted sexually, causing painful grouped vesicles/ulcers, urethritis, dysuria, and lymphadenopathy; primary infection may cause fever and systemic symptoms. Recurrent episodes are shorter and milder. Complications: aseptic meningitis, neonatal herpes (serious—transmission during delivery, prevention by suppressive therapy/cesarean), increased HIV risk. Diagnosis: PCR (preferred) or viral culture of lesions; type-specific serology. Management: antivirals (acyclovir, valacyclovir, famciclovir) for primary, episodic, or suppressive therapy (reduces recurrences and transmission); safe sex counseling; cesarean if active lesions at delivery.

\medterm{Genital Warts}-- Condyloma acuminata: benign epithelial proliferations caused by low-risk human papillomavirus (types 6, 11), transmitted sexually; pink or flesh-colored papules that may be flat, papular, or pedunculated, occurring on the genitalia, perianal region, and perineum. Usually asymptomatic but may cause itching, bleeding, or psychological distress. Diagnosis: clinical; biopsy if atypical. Management: patient-applied (podophyllotoxin, imiquimod, sinecatechins) or provider-administered (cryotherapy, TCA, electrocautery, laser, excision) treatments; recurrence is common; spontaneous regression possible. HPV vaccination prevents infection; treat sexual partners and screen for other STIs.

\medterm{Genital Warts (Condylomata Acuminata)}-- Anogenital warts caused by human papillomavirus (HPV), most commonly types 6 and 11 (low-risk, >90\% of genital warts, not associated with malignancy). Transmission: direct skin-to-skin contact during vaginal, anal, or oral sex; highly infectious (65\% transmission rate per sexual encounter). Incubation period: 2 weeks to 8 months (average 3 months). Clinical presentation: flesh-coloured, pink, or hyperpigmented papules or plaques on the genital, perianal, or oral mucosa. Morphologies: exophytic/cauliflower-like (condylomata acuminata), flat papules (smooth, slightly raised, more common on the cervix or penile shaft), keratotic (thick, wart-like, on dry skin of penile shaft or scrotum). Sites: males (foreskin, frenulum, glans, corona, shaft, scrotum, perianal), females (vulva, vaginal introitus, labia, perineum, perianal, cervix). Subclinical infection: detected only by HPV DNA testing or colposcopy. Natural history: 30\% regress spontaneously within 4 months, 70\% within 2 years. Malignant potential: HPV 6 and 11 rarely cause malignant transformation; HPV 16 and 18 (high-risk types) cause cervical, anal, and oropharyngeal cancer. Diagnosis: clinical (visual inspection, colposcopy, anoscopy). Biopsy if atypical, pigmented, ulcerated, or treatment-resistant. Treatment: patient-applied (imiquimod 5\% cream, podophyllotoxin 0.5\% solution/gel, sinceatechins 15\% ointment); provider-administered (cryotherapy, trichloroacetic acid 80-90\%, surgical excision, electrosurgery, laser ablation). No single treatment is universally effective; recurrence rates are 20-50\% regardless of modality. Prevention: HPV vaccine (quadrivalent Gardasil targets HPV 6, 11, 16, 18; nonavalent Gardasil 9 targets 9 types) prevents >90\% of genital warts when given before sexual debut.

\medterm{German Measles (Rubella)}-- See Rubella.

\medterm{Giardiasis}-- See gastroenterology.

\medterm{Glandular Fever (Infectious Mononucleosis)}-- See Mononucleosis.

\medterm{Gonorrhea}-- A sexually transmitted infection caused by the Gram-negative diplococcus Neisseria gonorrhoeae. Epidemiology: second most common bacterial STI (after Chlamydia), highest incidence in 15-24 year olds. Clinical features: (1) Men: acute urethritis (purulent discharge, dysuria, 2-7 day incubation); (2) Women: often asymptomatic (50-70\%), may present with cervicitis (mucopurulent discharge, intermenstrual bleeding, dyspareunia), urethritis, PID (salpingitis, tubo-ovarian abscess); (3) Extra-genital: pharyngeal (oropharyngeal swab--often asymptomatic), rectal (proctitis, discharge, tenesmus); (4) Disseminated: arthritis-dermatitis syndrome (tenosynovitis, migratory polyarthritis, pustular rash), perihepatitis (Fitz-Hugh-Curtis syndrome), meningitis, endocarditis (rare). Neonates: ophthalmia neonatorum (purulent conjunctivitis, can cause blindness without prophylaxis). Diagnosis: nucleic acid amplification test (NAAT) on urine or swab, culture for antibiotic sensitivity. Complications: PID, ectopic pregnancy, infertility (tubal scarring), epididymitis. Treatment: ceftriaxone 500 mg IM (single dose) plus azithromycin 1 g orally (single dose). Antimicrobial resistance: resistance to penicillin, tetracyclines, fluoroquinolones now widespread. Requires partner notification and treatment.

\medterm{Grippe (see Influenza)}-- An older term for influenza. See Influenza.

\medterm{Group A Strep}-- Group A Streptococcus (Streptococcus pyogenes), a gram-positive coccus in chains, beta-hemolytic. Infections: pharyngitis/tonsillitis, impetigo, erysipelas, cellulitis, scarlet fever, necrotizing fasciitis, streptococcal toxic shock syndrome, pneumonia, and puerperal sepsis. Complications: acute rheumatic fever (carditis, arthritis, chorea—post-pharyngitis) and post-streptococcal glomerulonephritis (post-skin/pharynx). Diagnosis: rapid antigen test/throat culture (pharyngitis), culture. Management: penicillin (or amoxicillin) is first-line for pharyngitis—prevents rheumatic fever; surgical debridement for necrotizing fasciitis; penicillin for invasive disease with clindamycin; treatment of close contacts for severe disease.

\medterm{Hansen's Disease (Leprosy)}-- A chronic infectious disease caused by Mycobacterium leprae, an acid-fast, obligate intracellular bacillus. Transmission: respiratory droplets (prolonged close contact). Incubation: 2-10 years (average 5). Epidemiology: 200,000 new cases per year globally (mostly India, Brazil, Indonesia). Pathophysiology: M. leprae has a predilection for skin and Schwann cells of peripheral nerves; the clinical spectrum depends on host cell-mediated immunity. Ridley-Jopling classification: (1) Tuberculoid leprosy (TT): strong cell-mediated immunity, few well-defined anesthetic skin lesions (hypopigmented macules with raised edges), thickened nerves, few bacilli. (2) Lepromatous leprosy (LL): poor cell-mediated immunity, numerous symmetric skin lesions (nodules, plaques, diffuse infiltration), extensive nerve involvement (symmetrical, stocking-glove anesthesia), many bacilli (positive slit-skin smear). (3) Borderline forms (BT, BB, BL): unstable, may shift towards TT or LL. Reactions: type 1 (reversal reaction, cell-mediated, worsening skin and nerve lesions), type 2 (erythema nodosum leprosum, immune complex-mediated, fever, painful nodules, neuritis, iritis). Complications: peripheral neuropathy (anesthesia leading to unrecognised injuries, neuropathic ulcers, Charcot joints), claw hand, foot drop, lagophthalmos (facial nerve), nasal deformity, blindness. Diagnosis: clinical (anesthetic skin lesions, thickened nerves), slit-skin smear for acid-fast bacilli, skin or nerve biopsy. Treatment: multidrug therapy (MDT): dapsone + rifampicin +/- clozimine, duration 6-12 months for paucibacillary, 12-24 months for multibacillary. Free MDT from WHO.

\medterm{Hantavirus}-- A genus of enveloped negative-sense RNA viruses of the Bunyaviridae family, transmitted to humans primarily through inhalation of aerosolized excreta.

\medterm{Healthcare-Associated Pneumonia}-- HCAP concept retired (2016 IDSA/ATS); risk factors for MDR now assessed individually. Risk fac-.

\medterm{Hepatitis}-- Inflammation of the liver caused by hepatotropic viruses. See Hepatitis (Viral Hepatitis) for complete definition.

\medterm{Hepatitis A}-- An acute viral infection of the liver caused by hepatitis A virus (HAV), transmitted via fecal-oral route (contaminated food or water). Incubation 15-50 days. Presents with fever, malaise, nausea, jaundice, and elevated transaminases. Self-limited; does not cause chronic infection. Fulminant hepatitis is rare. Diagnosis: IgM anti-HAV. Prevention: HAV vaccine (recommended for travelers, children) and hygiene. Treatment: supportive care.

\medterm{Hepatitis A, B, C, D, E}-- The five hepatotropic viruses that cause inflammation of the liver. See individual entries: Hepatitis A, Hepatitis B, Hepatitis C, Hepatitis D, Hepatitis E.

\medterm{Hepatitis B}-- A viral infection caused by hepatitis B virus (HBV), transmitted via blood, sexual contact, and perinatally. Can cause acute (self-limited or fulminant) and chronic infection (defined by HBsAg positivity >6 months). Chronic HBV leads to cirrhosis and hepatocellular carcinoma. Diagnosis: HBsAg, anti-HBs, HBeAg, HBV DNA, and liver function tests. Treatment: tenofovir, entecavir, or peg-interferon for chronic infection. Prevention: HBV vaccine (universal infant immunization).

\medterm{Hepatitis C}-- A viral infection caused by hepatitis C virus (HCV), primarily transmitted through blood exposure (IV drug use, blood transfusion before 1992). Acute infection is often asymptomatic; 50-80\% progress to chronic hepatitis. Chronic HCV can lead to cirrhosis and hepatocellular carcinoma over decades. Diagnosis: anti-HCV antibody, HCV RNA. Treatment: direct-acting antivirals (DAAs) achieve cure rates >95\%. No vaccine available. Screening recommended for at-risk populations.

\medterm{Hepatitis (Viral Hepatitis)}-- Inflammation of the liver caused by hepatotropic viruses. Hepatitis A (HAV): fecal-oral, acute self-limited, no chronic carrier state; vaccine-preventable. Hepatitis B (HBV): blood/sexual/perinatal transmission; acute and chronic infection; leading cause of cirrhosis and HCC. Hepatitis C (HCV): blood-borne; chronic in 50-80\%, leading cause of liver transplantation. Hepatitis D (HDV): requires HBV co-infection. Hepatitis E (HEV): fecal-oral, self-limited but severe in pregnancy. Diagnosis: viral serology and PCR. Treatment: supportive for HAV/HEV; antivirals for HBV (tenofovir, entecavir) and HCV (DAAs).

\medterm{Hepatitis Viruses}-- Hepatitis A: fecal-oral, acute self-limited, vaccine (inactivated). Hepatitis B: blood/sexual/perinatal. Acute: HBsAg, IgM anti-HBc. Chronic: HBsAg >6 months, HBV DNA, HBeAg. Treatment: entecavir, tenofovir (TAF/TDF) - suppress DNA, prevent cirrhosis/HCC. HDV: requires HBV, pegylated interferon. Hepatitis C: blood-borne. Acute: HCV RNA. Chronic: genotype (1-6), fibrosis (F0-F4, elastography/biopsy). DAA: pangenotypic (glecaprevir/pibrentasvir 8w, sofosbuvir/velpatasvir 12w, sofosbuvir/velpatasvir/voxilaprevir for retreatment). SVR12 = cure. Hepatitis E: fecal-oral (genotype 1/2), zoonotic (3/4, pork), severe in pregnancy. Hepatitis G/GBV-C: non-pathogenic.

\medterm{Herpes}-- A common viral infection caused by herpes simplex virus (HSV, types 1 and 2) and varicella-zoster virus (VZV, herpes zoster/shingles). See Herpes Simplex and Herpes Zoster.

\medterm{Herpes Simplex}-- A viral infection caused by herpes simplex virus type 1 (HSV-1, primarily orolabial) and type 2 (HSV-2, primarily genital). Primary infection followed by lifelong latency in sensory ganglia with periodic reactivation. Oral herpes: cold sores/fever blisters. Genital herpes: painful vesicular lesions, recurrent outbreaks. Complications: herpes keratitis, encephalitis, and neonatal herpes. Treatment: acyclovir, valacyclovir, famciclovir for acute episodes and suppressive therapy.

\medterm{Herpes Simplex (Cold Sores / Genital Herpes)}-- A viral infection caused by herpes simplex virus type 1 (HSV-1, typically orofacial, though increasingly anogenital) and type 2 (HSV-2, predominantly anogenital). HSV is a double-stranded DNA virus of the Herpesviridae family with latency in sensory ganglia and periodic reactivation. Primary infection: often asymptomatic; when symptomatic, causes painful vesicular lesions on an erythematous base. Orofacial herpes (cold sores): HSV-1, vesicular lesions on the vermilion border of the lips (herpes labialis), triggered by sun exposure, fever, stress, immunosuppression. Genital herpes: HSV-2 (80\%) and HSV-1 (20\%), painful genital/perianal vesicles that ulcerate and crust over; primary episode lasts 2-4 weeks; recurrent episodes are milder (viral shedding for 3-5 days). Other manifestations: herpetic whitlow (finger), herpes gladiatorum (skin-to-skin contact sports), herpes keratitis (corneal infection, leading cause of infectious blindness in developed countries), eczema herpeticum (Kaposi varicelliform eruption in atopic dermatitis), herpes encephalitis (temporal lobe, most common cause of sporadic viral encephalitis in adults). Neonatal herpes: transmitted during vaginal delivery if active genital lesions, severe with high mortality. Diagnosis: clinical, viral PCR (golden standard), Tzanck smear (multinucleated giant cells). Treatment: acyclovir, valacyclovir, famciclovir. Antiviral therapy can reduce duration and severity of primary and recurrent episodes. Suppressive therapy (daily valacyclovir 500 mg) reduces recurrences by 70-80\% and reduces transmission risk. No cure exists.

\medterm{Herpes Zoster (Shingles)}-- A painful, vesicular rash caused by reactivation of latent varicella-zoster virus (VZV, human herpesvirus 3) in the dorsal root ganglia. VZV causes varicella (chickenpox) as a primary infection, then remains latent in sensory ganglia; reactivation occurs when cell-mediated immunity declines (age >50, immunosuppression, HIV, malignancy, chemotherapy, radiotherapy, stress). Incidence: 1 in 3 people develop shingles in their lifetime; incidence rises sharply after age 50. Prodrome: 2-3 days of burning, stabbing pain, paraesthesia in a dermatomal distribution before the rash appears. Rash: unilateral, vesicular, dermatomal (thoracic most common 50\%, trigeminal--V1 ophthalmic branch 15\%), clusters of clear vesicles on erythematous base, pustules, crusting over 7-10 days, typically resolves within 2-4 weeks. Complications: Postherpetic neuralgia (PHN, most common, 10-20\% of patients, defined as pain persisting >90 days after rash onset, incidence increases with age; treatment: gabapentin, pregabalin, tricyclic antidepressants, lidocaine patch, capsaicin cream), herpes zoster ophthalmicus (Hutchinson sign: rash on tip of nose indicates nasociliary nerve involvement, increased risk of keratitis, uveitis, glaucoma, blindness; needs urgent ophthalmology referral), Ramsay Hunt syndrome (VZV in geniculate ganglion: ear pain, vesicles in auditory canal and pinna, facial palsy, tinnitus, hearing loss), disseminated zoster (visceral involvement in severely immunocompromised patients), motor neuropathy (rare, segmental weakness). Diagnosis: clinical, PCR or DFA of vesicle fluid. Treatment: antiviral therapy (acyclovir 800 mg 5x/day, valacyclovir 1 g TID, famciclovir 500 mg TID for 7 days) best within 72 hours of rash onset, reduces acute pain and PHN risk; analgesics (NSAIDs, opioids, gabapentinoids). Prevention: recombinant zoster vaccine (Shingrix, adjuvanted subunit glycoprotein E vaccine, 97\% efficacy in adults >50, 2 doses 2-6 months apart, recommended for all immunocompetent adults >50).

\medterm{Histoplasmosis}-- A systemic fungal infection caused by Histoplasma capsulatum, a dimorphic fungus found in nitrogenrich soil enriched with bat or bird guano, most commonly in the Ohio and Mississippi River valleys of the United States, Central and South America, and parts of Africa and Asia. The organism exists as a mold in the environment (producing infectious microconidia that are aerosolized when soil is disturbed) and converts to a yeast form at body temperature (37 degrees C). Infection occurs through inhalation of conidia, which reach the alveoli and are ingested by alveolar macrophages, where they convert to yeasts and multiply intracellularly. Most immunocompetent individuals (90 percent) develop asymptomatic or self-limited, mild influenza-like illness. The remaining 10 percent develop progressive forms depending on the intensity of exposure and host immune status. Acute pulmonary histo-.

\medterm{HIV}-- Human immunodeficiency virus: a retrovirus (HIV-1, HIV-2) causing progressive CD4+ T-cell depletion and AIDS if untreated. Transmitted via sexual contact, blood/blood products, mother-to-child. Natural history: acute seroconversion illness (fever, rash, lymphadenopathy), asymptomatic phase, symptomatic phase (opportunistic infections), AIDS (CD4 <200 or AIDS-defining illness). Diagnosis: HIV antibody/antigen test (4th generation), confirmatory viral load; seroconversion window. Management: ART (combination antiretroviral therapy) suppresses viral load, restores immunity, prevents transmission (U=U); pre-exposure (PrEP) and post-exposure (PEP) prophylaxis; monitor CD4 and viral load; treat/prophylax opportunistic infections.

\medterm{HIV/AIDS}-- HIV-1/2: retrovirus, CD4+ T cell depletion. Transmission: sexual, blood, perinatal, breast milk. Acute retroviral syndrome: fever, lymphadenopathy, pharyngitis, rash, high viral load. Clinical latency: CD4 decline over years. AIDS: CD4 <200 or AIDS-defining illness (PJP, toxo, CMV, MAC, crypto, KS, lymphoma, TB). ART: 3+ drugs from 2 classes (2 NRTI + INSTI/NNRTI/PI). Preferred: bictegravir/TAF/FTC, dolutegravir/abacavir/lamivudine (HLA-B*5701 negative), dolutegravir + TAF/FTC. Monitoring: viral load (target <50 copies/mL), CD4, resistance if failure. OI prophylaxis: PJP (TMPSMX if CD4<200), Toxo (if CD4<100), MAC (azithromycin if CD4<50), fungal (fluconazole if high risk). IRIS: paradoxical worsening after ART initiation. PrEP: TDF/FTC or TAF/FTC daily, cabotegravir LA bimonthly. PEP: 3-drug regimen within 72h x28d.

\medterm{HIV (Human Immunodeficiency Virus)}-- A retrovirus (family Retroviridae, genus Lentivirus) that causes acquired immunodeficiency syndrome (AIDS) by infecting and destroying CD4+ T lymphocytes (helper T cells). Two types: HIV-1 (pandemic, most common worldwide) and HIV-2 (less virulent, West Africa). HIV is an enveloped RNA virus; its genome encodes structural proteins (Gag, Env), enzymes (Pol--reverse transcriptase, integrase, protease), and regulatory proteins (Tat, Rev, Nef, Vif, Vpr, Vpu). Transmission: sexual contact (most common worldwide), blood-borne (IV drug use, blood transfusion), mother-to-child (vertical, during pregnancy, labour, or breastfeeding). Pathophysiology: HIV enters CD4+ cells via binding of gp120 to CD4 receptor and co-receptors (CCR5 or CXCR4); reverse transcriptase converts RNA to DNA; integrase incorporates viral DNA into host genome; protease cleaves viral polyproteins for assembly. Without treatment: acute HIV syndrome (2-4 weeks after infection: fever, rash, lymphadenopathy, pharyngitis, myalgia), clinical latency (8-12 years, viral replication continues, CD4 declines), AIDS (CD4 <200 cells/microL, opportunistic infections--PCP, cerebral toxoplasmosis, CMV retinitis, MAC, cryptococcal meningitis, and AIDS-defining cancers--Kaposi sarcoma, non-Hodgkin lymphoma, cervical cancer). Diagnosis: 4th generation antigen/antibody combo test (detects p24 antigen and HIV antibodies), confirm with Western blot or HIV RNA PCR. Management: antiretroviral therapy (ART) for life--combination therapy with 2 NRTIs (tenofovir disoproxil fumarate/emtricitabine) plus a third agent (integrase inhibitor--dolutegravir, bictegravir; NNRTI; or boosted PI). ART reduces viral load to undetectable (<20 copies/mL), restores immune function, prevents AIDS progression, and eliminates sexual transmission risk (U=U--undetectable equals untransmittable). Pre-exposure prophylaxis (PrEP: tenofovir/emtricitabine daily or on-demand for HIV-negative individuals at high risk). Post-exposure prophylaxis (PEP: 28 days of ART, started within 72 hours of exposure). Prevention: condoms, needle exchange, circumcision (reduces female-to-male transmission by 60\%), prevention of mother-to-child transmission (ART during pregnancy + formula feeding). Global: 39 million people living with HIV; 1.3 million new infections/year; 630,000 deaths/year.

\medterm{Hookworm}-- Intestinal nematode infection by Ancylostoma duodenale or Necator americanus, acquired by skin penetration of filariform larvae (walking barefoot on contaminated soil). Larvae migrate through lungs, adults attach to small intestine mucosa and suck blood, causing iron-deficiency anemia, protein loss, and malnutrition; ground itch (entry site). Diagnosis: stool microscopy for eggs. Management: albendazole or mebendazole; iron supplementation for anemia; prevention: footwear, sanitation, and mass drug administration programs. Chronic infection in children impairs growth and cognitive development.

\medterm{Host}-- A person or other living organism that can be infected by a virus or other pathogen under natural.

\medterm{HPV}-- Human papillomavirus: a DNA virus with >100 types causing epithelial infections. Low-risk types (6, 11) cause genital warts (condyloma acuminata) and laryngeal papillomas; high-risk types (16, 18—70 percent of cervical cancer; 31, 33, 45, etc.) cause cervical, anal, penile, vulvar, vaginal, oropharyngeal cancers and precancerous lesions. Transmission: sexual and skin contact; most infections clear spontaneously. Screening: cervical cytology (Pap) ± HPV DNA testing. Prevention: HPV vaccine (9-valent, 9-45 years—prevents warts and high-risk types). Management: treatment of warts (topical, cryotherapy, excision), treatment of CIN/precancer (LECP, conization), and surveillance of cervical abnormalities.

\medterm{Hydrophobia}-- See Rabies.

\medterm{Immunization}-- Routine childhood (CDC schedule): HepB, RV, DTaP, Hib, PCV, IPV, Influenza, MMR, VAR, HepA, HPV, MenACWY, MenB, COVID19. Adult: Tdap/Td q10y, influenza annually, zoster (RZV 2-dose), pneumococcal (PCV20 or PCV15+PPSV23), HPV (through 26, shared decision 27-45), HepA/B, MenACWY/MenB (riskbased), COVID-19. Pregnancy: Tdap 27-36w, influenza, COVID-19, RSV (32-36w). Contraindications: anaphylaxis to component, live vaccines in immunocompromised/pregnancy (MMR, VAR, LAIV, yellow fever, oral typhoid). Precautions: moderate/severe illness, GBS <6w post-vaccine. Vaccine safety: VAERS, VSD, CISA. Hesitancy: communication (presumptive recommendation, motivational interviewing).

\medterm{Incubation period}-- The time between when a person is exposed to an infection and when symptoms appear.

\medterm{Infection}-- See general.

\medterm{Infections in Pregnancy}-- Listeriosis: Listeria monocytogenes (deli meats, soft cheese), meningoencephalitis, fetal loss, neonatal sepsis. Ampicillin + gentamicin. GBS: vaginal colonization 20-30.

\medterm{Infective Endocarditis}-- See cardiology.

\medterm{Intravascular Catheter Infections}-- CLABSI: central line-associated bloodstream infection. Diagnosis: quantitative blood cultures (catheter vs peripheral >5:1), differential time to positivity (>2h), catheter tip culture (>15 CFU). Catheter-related BSIs: coagulase-negative staphylococci (most common), S. aureus, enterococci, gram-negatives, Candida. Management: remove catheter if S. aureus, Pseudomonas, Candida, mycobacteria, tunnel infection, port abscess, severe sepsis, persistent bacteremia >72h. Exchange over guidewire if no signs of exit site/tunnel infection and non-severe pathogens (CoNS, Bacillus, Corynebacterium). Antibiotics: CoNS 5-7d, S. aureus 2w (4-6w if endocarditis), gram-negative 7-14d, Candida 14d post-clearance. Lock therapy: antibiotic/ethanol/heparin locks for salvage. Prevention: maximal sterile barrier, chlorhexidine prep, daily necessity review, chlorhexidine dressings.

\medterm{Keshan disease}-- Heart disease caused by a lack of selenium, an element that the body needs to function properly.

\medterm{Koplik Spots}-- Pathognomonic enanthem of Measles (Rubeola virus infection). Consists of small, 1-2 mm, bluish-white spots on an erythematous background located on the buccal mucosa opposite the lower molars ("grains of salt on a red background"). Appear 1-2 days prior to the onset of the maculopapular rash and disappear 1-2 days after rash onset. Accompanied by the classic 3 Cs: Cough, Coryza, and Conjunctivitis.

\medterm{Lassa Fever}-- An acute viral hemorrhagic illness caused by Lassa virus, an arenavirus, endemic in West Africa (Nigeria, Sierra Leone, Guinea, Liberia). Transmission: exposure to excreta (urine, feces) of the multimammate rat (Mastomys natalensis), the natural reservoir. Incubation: 6-21 days. Clinical features: 80\% of infections are asymptomatic or mild. Symptomatic illness: gradual onset of fever, malaise, headache, sore throat, myalgia, retrosternal pain, cough, and gastrointestinal symptoms (vomiting, diarrhea). Severe disease (20\%): hemorrhage (gingival, nasal, gastrointestinal, vaginal), facial/neck edema, pleural effusion, hypotension/shock, encephalopathy, seizures, multi-organ failure, and sensorineural hearing loss (a unique complication, occurring in 25-30\% of survivors, may be permanent). Case fatality rate: 1-2\% overall, 15-20\% in hospitalised patients, >50\% in third trimester of pregnancy (high maternal and fetal mortality). Diagnosis: RT-PCR (most sensitive, first week), Lassa virus-specific IgM/IgG (seroconversion by week 3), virus isolation (BSL-4). Treatment: ribavirin (IV or oral, most effective if given early, within 6 days of onset). Supportive care: fluid resuscitation, blood products, oxygen, vasopressors. Prevention: rodent control, safe food storage, avoidance of rodent contact. No licensed vaccine. Outbreak management: strict infection control (contact and droplet precautions), case surveillance, contact tracing.

\medterm{Legionnaires' Disease}-- A severe form of pneumonia caused by Legionella pneumophila, a Gram-negative aerobic bacillus that is found naturally in freshwater environments (lakes, streams) and artificially in water systems (air conditioning cooling towers, hot tubs, decorative fountains, plumbing systems, showers, humidifiers). First identified after an outbreak at the 1976 American Legion convention in Philadelphia. Transmission: inhalation of aerosolised contaminated water (not person-to-person). Incubation: 2-14 days. Risk factors: age >50, smoking, chronic lung disease, immunocompromised (especially organ transplant recipients), diabetes, and male sex. Clinical features: atypical pneumonia--fever (>39 C), cough (initially non-productive), myalgia, headache, confusion, diarrhea, relative bradycardia, and hyponatraemia. Pontiac fever: a milder, self-limiting flu-like illness without pneumonia caused by the same bacterium. Diagnosis: urinary antigen test (detects L. pneumophila serogroup 1, 70-90\% sensitive), sputum culture (requires selective media, BCYE agar), PCR on respiratory samples, direct fluorescent antibody (DFA) staining. Treatment: macrolides (azithromycin 500 mg daily, 7-14 days) or fluoroquinolones (levofloxacin 750 mg daily, 7-14 days) as first-line. Doxycycline is an alternative. Severe cases require IV therapy. Complications: respiratory failure, septic shock, empyema, lung abscess, acute kidney injury, pericarditis. Mortality: 10-15\% (higher in immunocompromised). Reportable disease in most jurisdictions.

\medterm{Leishmaniasis}-- A parasitic disease caused by protozoa of the genus Leishmania, transmitted by the bite of infected female phlebotomine sandflies. Visceral leishmaniasis (kala-azar): caused by L. donovani/L. infantum -- fever, hepatosplenomegaly, pancytopenia, weight loss, hypergammaglobulinaemia; fatal if untreated. Cutaneous leishmaniasis: skin ulcers at bite sites that heal slowly with scarring (L. major, L. tropica, L. mexicana). Mucocutaneous leishmaniasis: destructive lesions of the nasopharynx (L. braziliensis). Diagnosis: microscopy of tissue smears (amastigotes in macrophages), culture, PCR. Treatment: liposomal amphotericin B (visceral), miltefosine, sodium stibogluconate.

\medterm{Leprosy}-- A chronic bacterial infection caused by Mycobacterium leprae, affecting skin and peripheral nerves. Transmission via nasal droplets. Spectrum of disease determined by immune response: tuberculoid (paucibacillary): one or few hypopigmented, anesthetic skin lesions, thickened nerves. Lepromatous (multibacillary): numerous symmetric nodules, plaques, and nerve involvement with sensory loss and deformities (leonine facies, claw hand, foot drop). Diagnosis: skin biopsy with acid-fast bacilli. Treatment: multi-drug therapy (dapsone + rifampin for paucibacillary; add clofazimine for multibacillary) for 6-12 months. Reactions require corticosteroids.

\medterm{Leprosy (Hansen's Disease)}-- See Hansen's Disease (Leprosy).

\medterm{Leptospirosis}-- A zoonotic bacterial infection caused by spirochaetes of the genus Leptospira, transmitted through contact with water or soil contaminated with the urine of infected animals (rodents, cattle, dogs). Incubation period 5-14 days. Mild form: fever, headache, myalgia (especially calves), conjunctival suffusion. Severe form (Weil disease): jaundice, renal failure, hemorrhage, myocarditis, and pulmonary hemorrhage with respiratory failure. Diagnosis: culture (blood/urine), PCR, serology (microscopic agglutination test). Treatment: doxycycline or penicillin for mild disease; intravenous penicillin or ceftriaxone for severe disease. Prevention: avoiding contaminated water, rodent control, doxycycline prophylaxis for high-risk exposure.

\medterm{Lice}-- Pediculosis: infestation by blood-sucking lice—head lice (Pediculus humanus capitis, common in children, scalp), body lice (Pediculus humanus corporis, clothing/seams, associated with poor hygiene and epidemic typhus transmission), and pubic/crab lice (Pthirus pubis, sexually transmitted). Features: pruritus (worse at night), visible nits (eggs) attached to hair shafts, excoriations; body lice with rash. Diagnosis: visual inspection, finding live lice/nits. Management: pediculicides (permethrin, malathion, benzyl alcohol, ivermectin), nit combing, laundering clothing/bedding, treating contacts; body lice—hygiene and laundering; treat with oral ivermectin for extensive.

\medterm{Listeriosis}-- A bacterial infection caused by Listeria monocytogenes, transmitted through contaminated food (deli meats, soft cheeses, unpasteurized dairy, produce). Risk groups: pregnant women (fever, flu-like illness), neonates (sepsis, meningitis), immunocompromised, and elderly. Manifestations: febrile gastroenteritis in immunocompetent; bacteremia and meningitis in at-risk. Neonatal: early-onset (sepsis, granulomatosis infantiseptica) and late-onset (meningitis) listeriosis. Diagnosis: blood and CSF culture. Treatment: ampicillin + gentamicin (synergistic); TMP-SMX for PCN-allergic.

\medterm{Lockjaw (Tetanus)}-- See Tetanus.

\medterm{Lupus Vulgaris}-- A chronic, progressive form of cutaneous tuberculosis (TB of the skin) caused by Mycobacterium tuberculosis. Once common in Europe, now rare in developed countries but still seen in developing regions. Presents as red-brown, gelatinous plaques (apple-jelly nodules on diascopy) typically on the face (nose, cheeks), which slowly enlarge and ulcerate over years without treatment. Can cause extensive tissue destruction and disfigurement. Diagnosis: skin biopsy (epithelioid granulomas with caseous necrosis), acid-fast stain, mycobacterial culture, PCR. Treatment: standard anti-TB therapy (isoniazid, rifampicin, ethambutol, pyrazinamide for 2 months, then isoniazid + rifampicin for 4 months).

\medterm{Lyme Disease}-- A tick-borne infection caused by Borrelia burgdorferi (Ixodes scapularis, deer tick), endemic to the northeastern and upper midwestern US. Early localized: erythema migrans (bullseye rash, >5cm) at tick bite site, fatigue, fever, myalgia. Early disseminated: multiple EM lesions, cranial nerve VII palsy (facial nerve), meningitis, carditis (AV block). Late: arthritis (monoarticular, knee), acrodermatitis chronica atrophicans, encephalopathy. Diagnosis: clinical (EM rash is diagnostic), two-tier serology (ELISA + Western blot). Treatment: doxycycline for early; ceftriaxone IV for neurologic or cardiac involvement. Prevention: tick avoidance, prophylactic doxycycline after high-risk tick bite.

\medterm{Lymphangitis}-- Inflammation of the lymphatic channels, most commonly due to acute bacterial infection spreading from a peripheral wound or cellulitis. The most common causative organism is group A Streptococcus pyogenes (90\%), followed by Staphylococcus aureus. Lymphangitis is a sign that infection is spreading and requires urgent medical attention. Clinical features: visible red, warm, tender linear streaks (lymphangitic streaking) extending proximally from the infected wound towards the draining lymph nodes (e.g., from the hand to the axilla, from the foot to the groin). Associated with: cellulitis, ascending lymphadenitis, fever, chills, malaise, and leukocytosis. The draining lymph nodes may be enlarged and tender (lymphadenitis). Diagnosis: clinical (pathognomonic appearance). Blood cultures may be positive. Complications: sepsis, bacteraemia, necrotising fasciitis (if anaerobes or mixed infection), and abscess formation. Treatment: antibiotics active against Streptococcus and Staphylococcus (e.g., cephalexin, clindamycin, flucloxacillin, or IV ceftriaxone/vancomycin for severe cases), elevation of the affected limb, warm compresses, analgesics. Surgical drainage if abscess forms. Prevention: prompt wound cleaning and antisepsis.

\medterm{Malaria}-- A life-threatening mosquito-borne disease caused by Plasmodium parasites (P. falciparum, P. vivax, P. ovale, P. malariae, P. knowlesi), transmitted by Anopheles mosquitoes. Presents with cyclical fevers, chills, headache, myalgia, and anemia. Severe P. falciparum malaria: cerebral malaria, severe anemia, respiratory distress, and multi-organ failure. Diagnosis: blood smear (thick and thin), rapid diagnostic test. Treatment: artemisinin-based combination therapy (ACT); severe malaria requires IV artesunate. Prevention: insecticide-treated nets, chemoprophylaxis.

\medterm{Meningococcal Disease}-- A bacterial infection caused by Neisseria meningitidis (serogroups A, B, C, W, Y), transmitted via respiratory droplets. Meningococcemia: sudden onset fever, petechial/purpuric rash, hypotension, and DIC. Meningococcal meningitis: fever, nuchal rigidity, headache, photophobia. Waterhouse-Friderichsen syndrome: fulminant meningococcemia with bilateral adrenal hemorrhage and septic shock. Diagnosis: blood and CSF culture, Gram stain, PCR. Treatment: IV ceftriaxone urgently; droplet precautions. Chemoprophylaxis for close contacts (rifampin, ciprofloxacin, ceftriaxone). Vaccination: MenACWY and MenB vaccines are available.

\medterm{Mononucleosis}-- An acute viral syndrome caused by Epstein-Barr virus (EBV), primarily affecting adolescents and young adults. Classic triad: fever, pharyngitis/exudative tonsillitis, and cervical lymphadenopathy. Associated with splenomegaly, fatigue, and atypical lymphocytosis. Diagnosis: positive heterophile antibody (Monospot) and EBV serology. Self-limited, supportive care. Avoid contact sports for 3-4 weeks due to splenic rupture risk.

\medterm{Neonatal Infections}-- Early-onset sepsis (EOS, <72h): GBS, E. coli, Listeria. Risk factors: maternal GBS+, chorioamnionitis, PROM >18h. Kaiser sepsis calculator for management. Late-onset sepsis (LOS, >72h): CoNS, S. aureus, gram-negatives, Candida. NEC: prematurity, formula feeding, ischemic gut. Pneumatosis intestinalis, portal venous gas. Treatment: NPO, antibiotics (ampicillin + gentamicin/clindamycin + metronidazole), surgery for perforation. HSV: SEM (skin/eye/mouth), CNS, disseminated. Acyclovir 20mg/kg q8h IV x14d (CNS/disseminated)

\medterm{Nocardiosis}-- A bacterial infection caused by Nocardia species (N. asteroides, N. brasiliensis), found in soil. Most common in immunocompromised patients (transplant, HIV, chronic steroid use). Pulmonary nocardiosis (most common): cough, fever, necrotizing pneumonia, cavitary lesions. Disseminated: brain abscesses (most common site), skin and soft tissue abscesses, bacteremia. Cutaneous nocardiosis: cellulitis, abscess, myectoma after trauma in immunocompetent. Diagnosis: modified acid-fast stain (partially acid-fast), culture (aerobic, grows slowly). Treatment: TMP-SMX (first-line); alternatives: carbapenems, linezolid, amikacin for severe or resistant cases. Prolonged therapy (6-12 months).

\medterm{Non-Tuberculous Mycobacteria}-- NTM: environmental mycobacteria (M. avium complex - MAC, M. kansasii, M. abscessus). Pulmonary NTM: ATS/IDSA criteria (clinical + radiographic + microbiologic: 2+ sputum or 1 BAL + nodules/cavities/bronchiectasis on CT). MAC: nodular bronchiectatic (middle-aged women, "Lady Windermere syndrome") or fibrocavitary (older men, COPD). Treatment: MAC macrolide (azithromycin/clarithromycin) + ethambutol + rifampin x12 months after culture conversion. M. kansasii - INH + RIF + EMB x12m post-conversion. M. abscessus - macrolide + IV amikacin + imipenem/cefoxitin + oral (clarithromycin) - difficult, often requires surgery. Disseminated MAC: AIDS (CD4 <50), MAC bacteremia, fever, weight loss, anemia. Treatment: clarithromycin/azithromycin + ethambutol + rifabutin; prophylaxis if CD4 <50 (azithromycin 1200mg weekly).

\medterm{Onchocerciasis}-- Also known as river blindness, a parasitic disease caused by the filarial nematode Onchocerca volvulus, transmitted by repeated bites of infected blackflies (Simulium species) breeding near fast-flowing rivers. Adult worms form subcutaneous nodules (onchocercomata). Microfilariae migrate through skin and eyes, causing intense pruritus, dermatitis, depigmentation (leopard skin), lichenification, and, most importantly, ocular damage leading to blindness (keratitis, iridocyclitis, chorioretinitis, optic atrophy). Diagnosis: skin snip biopsy (microfilariae), slit-lamp examination of the cornea, PCR. Treatment: ivermectin (annual or semi-annual mass drug administration, kills microfilariae but not adults). No safe, effective macrofilaricide exists.

\medterm{Opportunistic Infection}-- An infection caused by organisms that ordinarily do not cause disease in immunocompetent hosts but cause significant illness when host defences are impaired. Common causative organisms: Pneumocystis jirovecii (PCP pneumonia in HIV with CD4 <200/\(\mu\)L), Toxoplasma gondii (cerebral toxoplasmosis—HIV CD4 <100/\(\mu\)L), CMV (retinitis, colitis—HIV CD4 <50/\(\mu\)L, post-transplant), Mycobacterium avium complex (disseminated—HIV CD4 <50/\(\mu\)L), Cryptococcus neoformans (meningitis—HIV CD4 <100/\(\mu\)L), Aspergillus fumigatus (invasive aspergillosis—neutropenia, post-transplant), and Candida species (mucocutaneous and invasive—HIV, ICU, neutropenia). Prophylaxis is given based on CD4 thresholds in HIV (cotrimoxazole for PCP/Toxoplasma, azithromycin for MAC, fluconazole for Cryptococcus in endemic areas). Immune reconstitution inflammatory syndrome (IRIS) can occur when antiretroviral therapy is started in severely immunocompromised HIV patients.

\medterm{Orchitis}-- Inflammation of the testis, often associated with epididymitis (epididymo-orchitis). Causes: infectious—viral (mumps—most common cause of viral orchitis in postpubertal males, typically 4--6 days after parotitis, bilateral in 30\%), bacterial (Escherichia coli, Pseudomonas, sexually transmitted—Chlamydia trachomatis, Neisseria gonorrhoeae—typically ascending from urethra/epididymis), Coxsackie virus, and rarely tuberculosis. Non-infectious: trauma, autoimmune, and idiopathic. Presentation: acute scrotal pain and swelling, erythema, fever, and nausea. Diagnosis: clinical plus urine dipstick, urine culture, and Doppler ultrasound (increased testicular blood flow distinguishes orchitis from testicular torsion—torsion shows decreased flow). Treatment: antibiotics (ceftriaxone 500 mg IM single dose plus doxycycline 100 mg BD for 14 days for STI; levofloxacin for enteric organisms), NSAIDs, scrotal elevation, and cold packs; mumps orchitis is supportive (analgesia, steroids for severe inflammation). Complications: abscess, testicular infarction, atrophy, and infertility.

\medterm{Orf}-- A zoonotic viral skin infection caused by the parapoxvirus, transmitted to humans from infected sheep and goats (direct contact with lesions, fomites—virus survives in scabs for months). Orf is an occupational disease of farmers, butchers, and veterinarians. Incubation: 3--7 days. Lesions typically appear on the hands, fingers, or forearms: a single papule at the inoculation site progresses through stages—macule, target (bullseye) lesion, nodule (exudative, weeping), regenerative (dry crust), papillomatous, and regression (resolves over 4--6 weeks without scarring). Lesions may be large (up to 3--5 cm), tender, and occasionally complicated by secondary bacterial infection, lymphangitis, or erythema multiforme. Diagnosis is clinical; electron microscopy of vesicle fluid shows characteristic ovoid virions. Treatment: supportive (hygiene, dressings), no specific antiviral; lesions heal spontaneously. Prevention: gloves during handling of animals with orf lesions, quarantine of infected animals.

\medterm{Osteomyelitis}-- An infection of bone and bone marrow, most commonly caused by bacterial pathogens, with Staphylococcus aureus responsible for the majority of cases. It is classified by etiology (hematogenous, contiguous focus, or direct inoculation), duration (acute less than 2 weeks, subacute 2 weeks to 3 months, chronic greater than 3 months), and route. Hematogenous osteomyelitis typically affects the metaphysis of long bones in children (most commonly the femur, tibia, and humerus) and the vertebral bodies in adults (vertebral osteomyelitis).

\medterm{Outbreak}-- Synonymous with epidemic. Sometimes used to refer to a local epidemic as compared to a larger, general epidemic.

\medterm{Pandemic}-- A disease outbreak affecting large populations or a whole region, country, or continent.

\medterm{Parasite}-- An organism that lives on or within a host organism, deriving nutrients at the host's expense. Medical parasites are classified into three main groups: (1) Protozoa—single-celled organisms (Plasmodium—malaria, Toxoplasma—toxoplasmosis, Giardia—giardiasis, Entamoeba histolytica—amoebic dysentery, Leishmania—leishmaniasis, Trypanosoma—African sleeping sickness, Chagas disease). (2) Helminths—multicellular worms (nematodes/roundworms—Ascaris, hookworm, filariae; cestodes/tapeworms—Taenia solium, Echinococcus; trematodes/flukes—Schistosoma, liver flukes). (3) Ectoparasites—arthropods that infest the skin (scabies mites, lice, fleas, ticks, bed bugs). Parasitic infections disproportionately affect tropical and subtropical regions, causing significant morbidity and mortality (malaria alone causes >600,000 deaths annually). Transmission: vector-borne (mosquitoes, sandflies, ticks), fecal–oral (contaminated water/food), skin penetration (hookworm, Schistosoma), and vertical (Toxoplasma, malaria). Diagnosis: microscopy (blood films, stool examination, tissue biopsy), serology, antigen detection, and PCR. Treatment: antiparasitic drugs (antimalarials, antihelminthics—albendazole, praziquantel, ivermectin, metronidazole). Prevention: vector control, sanitation, hygiene, and health education.

\medterm{Parasitic Infections}-- Malaria: Plasmodium (falciparum, vivax, ovale, malariae, knowlesi). Falciparum: cerebral malaria, severe anemia, ARDS, AKI, hypoglycemia. Diagnosis: thick/thin smear, RDT (HRP2, pLDH),.

\medterm{Pathogen}-- A tiny organism such as a virus, bacterium, or parasite that can invade the body and produce disease.

\medterm{Pediculus humanus capitis}-- A blood-sucking parasite commonly known as the louse (plural, lice) that can cause an itchy scalp; infestations are highly contagious and especially common in school-age children.

\medterm{Pinworm}-- Infection with Enterobius vermicularis, the most common helminth infection in developed countries, primarily in school-age children. Lifecycle: ingestion of eggs $\rightarrow$ larvae hatch in small intestine $\rightarrow$ adults colonize the cecum/appendix; female migrates to perianal area at night to deposit eggs. Presents with intense perianal itching (especially at night), restless sleep, and secondary bacterial infection from scratching. Diagnosis: scotch tape test (adhesive tape applied to perianal area in the morning, examined microscopically for eggs). Treatment: albendazole or mebendazole (single dose, repeated after 2 weeks). Treat all household contacts.

\medterm{Plague}-- A bacterial infection caused by Yersinia pestis, transmitted by fleas from rodents. Three forms: bubonic (most common, 80-90\%): painful, swollen lymph nodes (buboes) after flea bite, fever, chills. Septicemic: bacteremia without bubo, endotoxic shock, DIC, purpura. Pneumonic: primary (inhalation) or secondary (from bubonic spread); fulminant pneumonia with bloody sputum; highly contagious via respiratory droplets. Untreated bubonic mortality 50-60\%; pneumonic nearly 100\% untreated. Diagnosis: culture, PCR, Gram stain of bubo aspirate. Treatment: streptomycin, gentamicin, doxycycline, or fluoroquinolones. Isolation for pneumonic plague.

\medterm{Polio}-- Abbreviation for poliomyelitis, an acute viral infection caused by poliovirus, an enterovirus of the Picornaviridae family transmitted via the fecal-oral route, primarily affecting children under 5 years of age. The virus replicates in the oropharynx and intestinal mucosa, then invades regional lymph nodes and the bloodstream; in a small proportion of cases (less than 1 percent), it invades the central nervous system and destroys motor neurons in the anterior horn of the spinal cord and brainstem. Most infections (approximately 95 percent) are asymptomatic or cause only mild nonspecific illness (abortive polio), while non-paralytic polio presents with fever, headache, vomiting, and meningismus. Paralytic polio presents with asymmetric acute flaccid paralysis, most commonly affecting the lower extremities, with loss of deep tendon reflexes but preserved sensation; bulbar polio affects the cranial nerves and respiratory center, potentially causing respiratory failure and death. There is no specific antiviral treatment; management is supportive, including physical therapy and respiratory support. Prevention is highly effective through vaccination with inactivated polio vaccine (IPV, Salk vaccine, injected) or oral polio vaccine (OPV, Sabin vaccine, oral). Global polio eradication efforts have reduced cases by over 99 percent since 1988, with wild poliovirus remaining endemic only in Afghanistan and Pakistan.

\medterm{Poliomyelitis}-- An acute viral infection caused by poliovirus (enterovirus), transmitted fecal-orally. Most infections are asymptomatic (95\%). Non-paralytic polio: fever, headache, vomiting, and meningismus. Paralytic polio: asymmetric acute flaccid paralysis due to destruction of anterior horn cells; most commonly affects lower extremities. Respiratory paralysis can be fatal. Post-polio syndrome: progressive muscle weakness decades after initial infection. Diagnosis: stool and throat culture, CSF PCR. No antiviral treatment; supportive care. Prevention: IPV (inactivated polio vaccine). Global eradication is near-complete.

\medterm{Prophylaxis}-- See general.

\medterm{Prostatitis}-- An inflammation of the prostate gland, sometimes caused by a bacterial infection, which may result in painful or difficult urination.

\medterm{Psittacosis}-- A zoonotic infection caused by Chlamydia psittaci, transmitted from infected birds (parrots, pigeons, poultry) via inhalation of dried bird droppings/secretions. Features: abrupt fever, headache, myalgia, dry cough, atypical pneumonia; complications: hepatitis, myocarditis, meningitis. Diagnosis: exposure history, serology (complement fixation, MIF), PCR. Management: doxycycline (or tetracycline) for 10-14 days; azithromycin alternative; reportable disease. Prevention: precautions with birds, treatment of infected birds.

\medterm{Pus}-- A thick, yellow or green liquid that is composed of dead cells and bacteria, most often found at the site of a bacterial infection.

\medterm{Pyrexia}-- Body temperature >38°C (100.4°F), also known as fever. A regulated rise in body temperature mediated by pyrogens (exogenous: microbial products, endotoxins; endogenous: IL-1, IL-6, TNF, prostaglandin E$_2$) acting on the hypothalamic thermoregulatory center. Pyrexia is a protective host response (enhances immune function, inhibits pathogen growth), not a disease itself. Fever of unknown origin (FUO): temperature >38.3°C for >3 weeks without diagnosis despite 3 outpatient visits or 3 inpatient days. Patterns: continuous (typhoid, UTI), intermittent (malaria, sepsis, abscess), relapsing (borreliosis, lymphoma), and Pel-Ebstein (Hodgkin lymphoma). Management: treat underlying cause; antipyretics (paracetamol, NSAIDs) for patient comfort. Fever itself does not require treatment unless >40°C or in vulnerable patients (young children, pregnancy, cardiac disease).

\medterm{Q Fever}-- A zoonotic disease caused by Coxiella burnetii, an obligate intracellular Gram-negative bacterium. Q fever is transmitted to humans primarily by inhalation of contaminated aerosols from infected livestock (cattle, sheep, goats) and their products (placenta, birth fluids, urine, feces, milk). Incubation: 2--3 weeks. Acute Q fever: high fever (up to 40°C, may last 1--3 weeks), severe headache, myalgia, arthralgia, fatigue, and atypical pneumonia (cough, chest pain, abnormal chest X-ray). Hepatitis (granulomatous—doughnut granulomas on liver biopsy) is common. Chronic Q fever (<5\% of cases): endocarditis (culture-negative, affecting prosthetic and native valves), vascular graft infection, and osteomyelitis. Diagnosis: serology (phase II IgG and IgM—antibodies rise 2--4 weeks after infection; phase I IgG \(\ge\)1:800 suggests chronic infection), PCR on blood or tissue during acute illness. Treatment: acute—doxycycline 100 mg BD for 14 days (first-line, shortens fever duration). Chronic—doxycycline plus hydroxychloroquine for \(\ge\)18 months. Prevention: pasteurisation of milk, biosecurity measures on farms, and veterinary vaccination (phase I vaccine—Coxevac). Q fever is a notifiable disease in many countries. The Netherlands experienced the largest outbreak (2007--2010, >4,000 cases) linked to intensive goat farming.

\medterm{Quarantine}-- See general.

\medterm{Relapsing Fever}-- An acute febrile illness caused by Borrelia species, characterized by recurrent episodes of fever separated by afebrile intervals. Two forms: (1) Louse-borne relapsing fever (Borrelia recurrentis)—transmitted by the human body louse (Pediculus humanus corporis), associated with overcrowding, poverty, and refugee settings; epidemic potential. (2) Tick-borne relapsing fever (Borrelia duttonii, B. hermsii, B. parkeri)—transmitted by soft ticks (Ornithodoros), endemic in Africa, North America, and the Middle East; tends to have more relapses. Pathophysiology: Borrelia spirochetes multiply in the blood and cause high fever; the immune system clears the organisms by antibody production, but antigenic variation (variable major proteins—Vmp) allows a new variant to emerge and cause a relapse. Clinical: sudden onset of high fever (39--40°C), rigors, severe headache, myalgia, arthralgia, photophobia, and hepatosplenomegaly. The initial febrile episode lasts 3--6 days, followed by an afebrile period of 7--9 days, then 1--3 relapses (tick-borne) or a single relapse (louse-borne). A petechial rash may be present. Diagnosis: dark-field microscopy of blood smears (spirochetes visible during febrile episodes); PCR. Treatment: doxycycline 100 mg BD for 7--10 days (or tetracycline, erythromycin). Jarisch–Herxheimer reaction (fever, rigors, hypotension) occurs within 2--4 hours of starting antibiotics in 50--90\% of patients (especially louse-borne)—pre-treated with IV fluids and antipyretics; severe cases may require hydrocortisone.

\medterm{River Blindness}-- See Onchocerciasis.

\medterm{Roundworm}-- Ascariasis: infection by Ascaris lumbricoides, the most common human soil-transmitted helminth, acquired by ingesting eggs from contaminated food/soil. Life cycle: larvae migrate through lungs (pneumonitis, Loeffler syndrome) $\rightarrow$ adult worms in small intestine (malnutrition, obstruction, biliary colic, intestinal obstruction in heavy infection); worms passed in stool. Diagnosis: stool microscopy (eggs), adult worms. Management: albendazole or mebendazole (single dose); prevention: hygiene, sanitation, and treatment of infected. Distinguish from other roundworms (hookworm, Strongyloides, Trichuris).

\medterm{Salmonella}-- A genus of Gram-negative, facultatively anaerobic, rod-shaped bacteria in the Enterobacteriaceae family, a major cause of foodborne illness worldwide. Classification: Salmonella enterica subsp. enterica serovars (>2,500 serovars). Clinically important serovars: S. Typhi and S. Paratyphi A, B, C (cause enteric fever—typhoid and paratyphoid fever—systemic illness, humans are the only reservoir, transmitted by fecal–oral route through contaminated food/water), and non-typhoidal Salmonella (NTS): S. Enteritidis and S. Typhimurium (most common—cause self-limited gastroenteritis, transmitted from animals—poultry, eggs, meat, reptiles). Pathogenesis: bacteria invade intestinal epithelial cells (via type III secretion system), survive within macrophages, and cause inflammation. Clinical: NTS gastroenteritis—nausea, vomiting, non-bloody diarrhea (may become bloody), abdominal cramps, fever (self-limiting, 4--7 days). Enteric fever (typhoid)—incubation 7--14 days, prolonged high fever (stepwise rise), headache, malaise, rose spots on trunk, relative bradycardia, hepatosplenomegaly, constipation initially then diarrhea. Complications of typhoid: intestinal perforation and hemorrhage (most feared). Chronic carriers: S. Typhi persists in the gallbladder (gallstones) and is shed in stool (Typhoid Mary). Diagnosis: stool culture (gastroenteritis), blood culture (enteric fever—most sensitive in first week), bone marrow culture (gold standard), and Widal test (serology, limited utility). Treatment: NTS gastroenteritis—supportive care (antibiotics not recommended for mild cases; prolongs carriage); severe or immunocompromised—ceftriaxone or azithromycin. Enteric fever—ceftriaxone 2 g IV daily for 10--14 days, azithromycin, or ciprofloxacin (if sensitive). Multidrug-resistant (MDR) and extensively drug-resistant (XDR) S. Typhi are emerging (azithromycin and carbapenems for XDR). Prevention: hand hygiene, safe food/water, typhoid vaccines (Vi polysaccharide injectable, Ty21a oral live attenuated).

\medterm{Scabies}-- A contagious skin infestation caused by the mite Sarcoptes scabiei var. hominis, transmitted by prolonged skin-to-skin contact. The female mite burrows into the stratum corneum, laying eggs, causing intense pruritus (worse at night) and a papular erythematous rash. Typical distribution: interdigital finger webs, wrists, elbows, axillae, waist, genitalia (spares the face in adults). Burrows appear as greyish, serpiginous lines. Crusted (Norwegian) scabies occurs in immunocompromised patients. Diagnosis: clinical, confirmed by microscopy of scrapings from burrows (mite, eggs, feces). Treatment: permethrin 5\% cream (two applications, 1 week apart) or oral ivermectin. Treat all close contacts.

\medterm{Sepsis}-- See emergency\_medicine.

\medterm{Sepsis Syndrome}-- The spectrum of the host's dysregulated response to infection: sepsis (infection + organ dysfunction, SOFA increase >=2), septic shock (sepsis + hypotension requiring vasopressors with lactate >2 mmol/L despite adequate fluid), and severe organ dysfunction/multi-organ failure. Features: fever/hypothermia, tachycardia, tachypnea, hypotension, altered mental status, raised lactate, elevated inflammatory markers. Management (Surviving Sepsis): early recognition, cultures before antibiotics, broad-spectrum antibiotics within 1 hour, IV fluids (30 mL/kg), vasopressors (norepinephrine), source control, and lactate-guided resuscitation; monitor organ function in ICU. High mortality.

\medterm{Septicaemia}-- See Sepsis.

\medterm{Septicemia}-- A condition in which disease-causing organisms have spread to the bloodstream from an infection elsewhere in the body. Also known as blood poisoning.

\medterm{Sexually Transmitted Infections (STIs)}-- Sexually transmitted infections are infections passed from one person to another through sexual contact, including vaginal, anal, and oral sex. Common STIs include chlamydia, gonorrhea, syphilis, herpes simplex virus, human papillomavirus, and HIV. Many STIs can be asymptomatic, making regular screening essential for sexually active individuals; treatment varies by infection, with bacterial STIs being curable with antibiotics, while viral STIs are managed with antiviral medications to reduce symptoms and transmission risk.

\medterm{Shingles}-- Herpes zoster: reactivation of varicella-zoster virus (VZV) from dorsal root ganglia, causing a painful, unilateral vesicular rash in a dermatomal distribution ('belt-like'), often preceded by burning pain. Risk increases with age and immunosuppression. Complications: postherpetic neuralgia (persistent pain), ophthalmic zoster (eye involvement—Ramsay Hunt syndrome with facial palsy if geniculate), disseminated disease, bacterial superinfection. Diagnosis: clinical; PCR/DFA of lesions. Management: antivirals (acyclovir, valacyclovir, famciclovir) within 72 h reduce duration and neuralgia; analgesics including gabapentin/amitriptyline; zoster vaccine (Shingrix) for prevention. Contagious to VZV-naive contacts (causes chickenpox).

\medterm{Smallpox}-- A highly contagious and deadly disease caused by variola virus, eradicated globally by 1980 through vaccination. Characterized by febrile prodrome followed by vesicular/pustular rash (centrifugal distribution, all lesions at same stage). Last known natural case: 1977. Bioterrorism concern due to stored virus in WHO-approved laboratories. Diagnosis: PCR, electron microscopy. No proven treatment; supportive care. Vaccination within 3-4 days of exposure can prevent disease.

\medterm{Staph Infection}-- Infections caused by Staphylococcus aureus (and coagulase-negative staphylococci), a gram-positive coccus in clusters. Skin/soft tissue (boils, cellulitis, abscess, impetigo), bacteremia, infective endocarditis, osteomyelitis, septic arthritis, pneumonia (post-influenza), toxic shock syndrome (TSST superantigen), food poisoning (enterotoxin), and device infections. MRSA (methicillin-resistant) is a major concern—community and hospital associated. Diagnosis: culture and sensitivity. Management: incision and drainage of abscesses, antibiotics (flucloxacillin/cefazolin for MSSA; vancomycin/daptomycin/linezolid for MRSA), treat endocarditis with prolonged therapy and source control; decolonization strategies.

\medterm{Staphylococcal Infection}-- Infections caused by Staphylococcus aureus and other staphylococci. Manifestations range from superficial skin infections (folliculitis, furuncles, carbuncles, impetigo) to deep-seated infections (osteomyelitis, septic arthritis, endocarditis, pneumonia, bacteremia, sepsis). Methicillin-resistant S. aureus (MRSA) is a major healthcare and community-associated pathogen. Treatment: beta-lactams for MSSA, vancomycin for MRSA; incision and drainage for abscesses.

\medterm{Streptococcal Infection}-- Infections caused by Streptococcus pyogenes (Group A Strep, GAS). Common presentations: pharyngitis (strep throat), scarlet fever (pharyngitis with erythematous rash), impetigo, cellulitis, and necrotizing fasciitis. Immune-mediated complications: acute rheumatic fever (can cause rheumatic heart disease) and post-streptococcal glomerulonephritis. Diagnosis: rapid antigen test, throat culture. Treatment: penicillin or amoxicillin. Prevention of rheumatic fever requires adequate treatment of GAS pharyngitis.

\medterm{Strongyloidiasis}-- Infection with Strongyloides stercoralis, a unique helminth capable of autoinfection. Most infections asymptomatic or mild (abdominal pain, diarrhea, urticaria). Larvae migration causes larva currens (rapidly migrating serpiginous rash, perianal). Hyperinfection syndrome: life-threatening in immunocompromised (HIV/HTLV-1, corticosteroids, transplant), leading to massive larval dissemination with Gram-negative bacteremia, meningitis, and respiratory failure. Diagnosis: stool microscopy (Baermann technique, agar plate culture), serology. Treatment: ivermectin (first-line); albendazole as alternative. Screening and treatment before immunosuppression is critical.

\medterm{Swine Flu}-- See Influenza.

\medterm{Systemic Inflammatory Response Syndrome}-- SIRS: a systemic inflammatory response defined by two or more of: temperature >38C or <36C, heart rate >90, respiratory rate >20 or PaCO2 <32 mmHg, WBC >12,000 or <4,000 or >10 percent bands. May result from infection (sepsis), trauma, burns, pancreatitis, surgery, or other insults. Historically used to define sepsis; now largely superseded by sepsis-3 criteria (SOFA, qSOFA) which emphasize organ dysfunction. SIRS criteria are sensitive but nonspecific; still used in pediatric definitions and bedside screening. Management: identify and treat the trigger, support organ function.

\medterm{Taeniasis (Tapeworm Infection)}-- Intestinal infection caused by adult tapeworms of the genus Taenia (cestodes). Two important species: Taenia solium (pork tapeworm) and Taenia saginata (beef tapeworm). Transmission: ingestion of undercooked meat containing cysticerci (larval cysts). Adult tapeworms attach to the small intestinal wall via scolex (armed with hooks in T. solium, suckers only in T. saginata) and produce proglottids (segments) containing eggs. Clinical: often asymptomatic; may cause mild abdominal discomfort, nausea, weight loss, increased appetite, and passage of proglottids in stool. Diagnosis: stool microscopy (eggs, proglottids—T. solium eggs are indistinguishable from T. saginata; species identified by uterine branching in gravid proglottids), perianal tape test. Treatment: praziquantel (5--10 mg/kg single dose) or niclosamide. Complications: T. solium can cause neurocysticercosis (ingestion of eggs from autoinfection or contaminated food/water)—larval cysts in the brain cause seizures, hydrocephalus, and focal neurological deficits. Prevention: thorough cooking of meat, improved sanitation, and meat inspection.

\medterm{Tapeworm}-- Cestode infection: segmented flatworms (tapeworms) transmitted by ingestion of larvae (pork—Taenia solium; beef—Taenia saginata; fish—Diphyllobothrium latum, causing B12 deficiency) or eggs (Echinococcus—hydatid cysts; Hymenolepis). Intestinal tapeworm: abdominal discomfort, passage of proglottids in stool; T. solium eggs cause cysticercosis (neurocysticercosis—seizures, intracranial cysts). Diagnosis: stool ova/parasite examination (proglottids, eggs), serology, imaging for cysticercosis/hydatid. Management: praziquantel (or niclosamide) for intestinal forms; albendazole ± surgery for cysticercosis and hydatid; prevention: thorough cooking of meat/fish, sanitation.

\medterm{Tapeworm Infection}-- See Taeniasis (Tapeworm Infection).

\medterm{Tetanus}-- An acute neurological disease caused by tetanospasmin toxin from Clostridium tetani, introduced through wounds. Characterized by generalized rigidity and painful muscle spasms (trismus/lockjaw, risus sardonicus, opisthotonos). Autonomic instability can occur in severe cases. Diagnosis: clinical. Treatment: human tetanus immune globulin (HTIG), wound debridement, metronidazole, and supportive care (benzodiazepines, ICU). Prevention: Tdap/Td vaccination at recommended intervals.

\medterm{Threadworm (Enterobiasis)}-- A common intestinal parasitic infection caused by Enterobius vermicularis (pinworm). Prevalence: most common helminth infection in developed countries, primarily affecting children aged 5--10 years. Transmission: fecal–oral (ingestion of eggs from contaminated hands, food, bedding, or fomites). Eggs are deposited in the perianal region by the female worm at night, causing intense perianal itching (pruritus ani), especially at night. Scratching contaminates fingers and nails, perpetuating the cycle. Diagnosis: perianal tape test (apply transparent adhesive tape to the perianal area in the morning before bathing, examined under microscope for eggs). Stool microscopy has low sensitivity. Treatment: mebendazole 100 mg single dose (or albendazole 400 mg single dose), repeated after 2 weeks (to kill newly hatched worms). All household members should be treated simultaneously. Good hand hygiene, short nails, frequent washing of bedding and pyjamas, and avoidance of scratching. Reinfection is common. Complications: vulvovaginitis (in girls), occasional appendicitis (worm in the appendix).

\medterm{Tinea}-- Dermatophytosis: superficial fungal infection of keratinized tissue (skin, hair, nails) by dermatophytes (Trichophyton, Microsporum, Epidermophyton). Types by site: tinea corporis (body—ringworm), tinea pedis (athlete's foot), tinea cruris (jock itch), tinea capitis (scalp—most common in children), tinea unguium/onychomycosis (nails), tinea manuum, tinea faciei. Features: annular erythematous scaly plaque with central clearing and active border. Diagnosis: clinical, KOH preparation, culture, or PCR. Management: topical antifungals (terbinafine, clotrimazole, miconazole) for most; oral terbinafine/griseofulvin for tinea capitis and nail infection; hygiene measures to prevent spread.

\medterm{Toxic Shock Syndrome}-- An acute, life-threatening illness caused by bacterial exotoxins. Staphylococcal TSS: associated with tampon use, surgical wounds, and nasal packing; presents with fever, hypotension, diffuse rash, and multi-organ involvement. Streptococcal TSS: associated with invasive GAS infections. Diagnosis: CDC case definition. Treatment: source control, IV fluids, vasopressors, and antibiotics (clindamycin + a beta-lactam for Strep; vancomycin for MRSA). IVIG may be beneficial.

\medterm{Toxocariasis}-- Infection by dog roundworm (Toxocara canis) or cat (T. cati), acquired by ingesting infective eggs from soil contaminated by animal feces (pica in children a risk). Larvae do not mature in humans but migrate (visceral larva migrans—fever, hepatomegaly, eosinophilia, pulmonary symptoms) or cause ocular larva migrans (granuloma in the eye, vision loss, strabismus). Diagnosis: clinical, eosinophilia, ELISA serology (Toxocara antibodies); eye examination. Management: albendazole or mebendazole ± corticosteroids (ocular/visceral severe); prevention: deworm pets, hygiene, prevent pica.

\medterm{Toxoid vaccines}-- Vaccines that protect against harmful bacterial toxins. These vaccines contain toxins that have been detoxified and rendered inactive.

\medterm{Toxoplasmosis}-- Infection caused by the intracellular protozoan Toxoplasma gondii. Definitive host: cats (shed oocysts in feces). Humans are intermediate hosts, infected via ingestion of undercooked meat (containing tissue cysts), cat feces, or transplacentally. Immunocompetent: usually asymptomatic or mild flu-like illness with lymphadenopathy. Immunocompromised (HIV, transplant): toxoplasma encephalitis (brain abscesses, headache, confusion, seizures, focal neurology), myocarditis, pneumonitis. Congenital toxoplasmosis: infection in pregnancy can cause chorioretinitis, intracranial calcifications, hydrocephalus, and developmental delay. Diagnosis: serology (IgM, IgG, PCR), imaging (CT/MRI). Treatment: pyrimethamine + sulfadiazine + folinic acid.

\medterm{Trachoma}-- A chronic keratoconjunctivitis caused by Chlamydia trachomatis (serovars A-C), the leading infectious cause of blindness worldwide. Repeated infection causes conjunctival scarring, trichiasis (in-turned eyelashes), corneal abrasion, opacification, and blindness. Transmission: direct contact, flies, fomites. Active trachoma: follicular conjunctivitis, papillary hypertrophy. Cicatricial trachoma: scarring, entropion, trichiasis. Diagnosis: clinical WHO grading system. Treatment: SAFE strategy (Surgery for trichiasis, Antibiotics: azithromycin mass drug administration, Facial cleanliness, Environmental improvement). Global elimination program is ongoing.

\medterm{Travel Medicine}-- Pre-travel consultation: vaccines (routine, hepatitis A/B, typhoid, yellow fever, meningococcal, rabies pre-exposure, Japanese encephalitis, cholera), malaria prophylaxis (atovaquone-proguanil daily, doxycycline daily, mefloquine weekly, tafenoquine.

\medterm{Trichinosis}-- A parasitic infection caused by Trichinella spiralis, acquired by eating undercooked pork or wild game (bear, boar) containing encysted larvae. Larvae are released in the intestine (enteric phase: diarrhea, abdominal pain), then migrate to skeletal muscle (parenteral phase: myositis, periorbital edema, fever, eosinophilia). Classic presentation: periorbital edema + myalgia + eosinophilia after eating undercooked meat. Diagnosis: serology (ELISA, Western blot), muscle biopsy. Treatment: albendazole or mebendazole (enteric phase); corticosteroids for severe myositis.

\medterm{Trichomonas}-- A protozoan parasite, Trichomonas vaginalis, causing trichomoniasis—the most common non-viral sexually transmitted infection worldwide (156 million cases annually). T. vaginalis is a flagellated, motile protozoan that infects the urogenital tract. Transmission: sexual intercourse. Incubation: 4--28 days. Clinical: women—vaginal discharge (frothy, yellow-green, malodorous), vulvovaginal itching, dysuria, dyspareunia, punctate cervical haemorrhages ('strawberry cervix' on colposcopy). Men—often asymptomatic (70\%), may cause urethritis (thin, white discharge), balanitis, prostatitis. Complications: preterm birth, low birth weight, and increased HIV acquisition/transmission. Diagnosis: wet mount microscopy (motile trichomonads, sensitivity 60--70\%), culture, nucleic acid amplification test (NAAT—gold standard, sensitivity >95\%). Treatment: metronidazole 2 g PO single dose (or tinidazole). Treat sexual partners. Metronidazole resistance occurs in 5--10\% (treat with higher dose or tinidazole). Avoid alcohol during and for 24 hours after treatment (disulfiram-like reaction).

\medterm{Trichomoniasis}-- A sexually transmitted infection caused by the protozoan Trichomonas vaginalis. Often asymptomatic in men; women present with vaginal discharge (frothy, yellow-green, malodorous), itching, dysuria, and inflammation. Associated with preterm birth and increased HIV acquisition risk. Diagnosis: wet mount microscopy, culture, NAAT. Treatment: metronidazole or tinidazole (single dose) for patient and partner(s).

\medterm{Trypanosomiasis}-- Two forms: African trypanosomiasis (sleeping sickness, transmitted by tsetse fly, caused by Trypanosoma brucei gambiense -- chronic West African form, or T. b. rhodesiense -- acute East African form) and American trypanosomiasis (Chagas disease, transmitted by triatomine bugs, caused by Trypanosoma cruzi). Clinical features: African -- chancre at bite site, fever, lymphadenopathy (Winterbottom sign), headache, followed by meningoencephalitis (sleep disturbance, neurological deterioration, coma). Chagas disease: acute phase (fever, Romana sign -- unilateral periorbital edema), chronic phase (cardiomyopathy, megaoesophagus, megacolon decades later). Diagnosis: microscopy of blood/CSF, serology, PCR. Treatment: African -- pentamidine/suramin (early), melarsoprol/eflornithine (late). Chagas -- benznidazole, nifurtimox.

\medterm{Tuberculosis}-- TB: Mycobacterium tuberculosis, airborne transmission. Primary TB: Ghon focus + hilar lymphadenopathy (Ghon complex). Reactivation TB: apical/posterior upper lobe cavitary lesions. Clinical: cough >2 weeks, fever, night sweats, weight loss, hemoptysis. Diagnosis: AFB smear (3 sputums, sensitivity 50-60.

\medterm{Tuberculosis (TB)}-- A chronic infection caused by Mycobacterium tuberculosis, transmitted by airborne droplets; pulmonary TB is most common, extrapulmonary forms occur (lymph node, pleural, meningeal, miliary, skeletal, renal). Features: chronic cough (>3 weeks), hemoptysis, fever, night sweats, weight loss; Ghon complex, apical cavitation on imaging. Diagnosis: sputum smear/culture, nucleic acid amplification (GeneXpert), IGRA/TST (latent), CXR. Treatment: RIPE (rifampicin, isoniazid, pyrazinamide, ethambutol) 2 months + rifampicin/isoniazid 4 months (6-month total); MDR/XDR requires second-line agents. Latent TB: isoniazid or rifampicin prophylaxis. Directly observed therapy and infection control are key.

\medterm{Tularemia}-- A zoonotic bacterial infection caused by Francisella tularensis, transmitted through tick bites, handling infected animals (rabbits, rodents), contaminated water, or aerosolization. Ulceroglandular (most common): ulcer at inoculation site + tender lymphadenopathy. Other forms: glandular (lymphadenopathy only), oculoglandular (conjunctivitis + preauricular adenopathy), oropharyngeal (pharyngitis), typhoidal (fever without localizing signs), pneumonic (cough, chest pain). Diagnosis: serology, PCR, culture (biohazard). Treatment: streptomycin, gentamicin, doxycycline, or ciprofloxacin.

\medterm{Typhus}-- A group of acute febrile illnesses caused by Rickettsia species, transmitted by arthropod vectors. Three main forms: (1) Epidemic (louse-borne) typhus—Rickettsia prowazekii, transmitted by the human body louse (Pediculus humanus corporis). Associated with overcrowding, poverty, cold climate, and refugees ('war typhus'). Incubation 10--14 days. Clinical: sudden high fever, severe headache, myalgia, prostration, a macular rash starting in the axillae and spreading to the trunk (spares face, palms, soles), and meningismus. Mortality 10--60\% without treatment (highest in older, malnourished patients). Brill–Zinsser disease: reactivation of latent R. prowazekii years after primary infection, milder. (2) Endemic (murine/flea-borne) typhus—Rickettsia typhi, transmitted by rat fleas (Xenopsylla cheopis). Milder than epidemic typhus; rash is maculopapular, centripetal. (3) Scrub typhus (tsutsugamushi disease)—Orientia tsutsugamushi (formerly Rickettsia tsutsugamushi), transmitted by chigger mites (Leptotrombidium). Endemic in the Asia-Pacific region ('tsutsugamushi triangle'). An eschar (tache noire) at the bite site is characteristic. Diagnosis: serology (IFA—gold standard; Weil–Felix test is obsolete), PCR on blood or eschar biopsy. Treatment: doxycycline 100 mg BD for 7--14 days (first-line). Azithromycin and chloramphenicol are alternatives (pregnancy, doxycycline allergy). Prevention: louse control (epidemic typhus—delousing with insecticides, improved hygiene), flea control (endemic typhus—rodent control), mite avoidance (scrub typhus—insect repellent DEET, protective clothing). No vaccine is currently available for typhus.

\medterm{Undulant Fever}-- See Brucellosis.

\medterm{Urinary Tract Infection}-- Infection of the urinary tract, classified as lower (cystitis, urethritis) or upper (pyelonephritis). Most common pathogen: Escherichia coli (80\%). Cystitis: dysuria, urinary frequency, urgency, suprapubic pain. Pyelonephritis: fever, flank pain, nausea/vomiting, costovertebral angle tenderness. Complicated UTI: structural/functional abnormality, catheter, pregnancy, male sex, immunocompromised. Diagnosis: urinalysis with nitrites and leukocyte esterase, urine culture. Treatment: TMP-SMX, nitrofurantoin, fosfomycin for uncomplicated; fluoroquinolones or cephalosporins for complicated/pyelonephritis.

\medterm{Urinary Tract Infections}-- UTI: significant bacteriuria (>105 CFU/mL) with symptoms. Uncomplicated cystitis: women, dysuria, frequency, urgency, suprapubic pain. Nitrite/leukocyte esterase dipstick. Pathogens: E. coli (75-90.

\medterm{Urinary Tract Infections (UTIs)}-- Urinary tract infections are bacterial infections affecting any part of the urinary system, including the urethra, bladder, ureters, and kidneys. They are much more common in women due to the shorter urethra, and typical symptoms include urinary urgency, frequency, dysuria, and suprapubic pain, while upper tract infections may cause flank pain, fever, and systemic illness. Diagnosis is confirmed by urinalysis and culture, and treatment involves antibiotic therapy tailored to the causative organism, with a emphasis on adequate hydration and preventive measures.

\medterm{UTI}-- See Urinary Tract Infection.

\medterm{Varicella}-- Varicella-zoster virus (VZV) infection: primary infection causes chickenpox, characterized by a generalized vesicular rash (macules $\rightarrow$ papules $\rightarrow$ vesicles $\rightarrow$ crusts, 'dew drops on a rose petal'), fever, and malaise; highly contagious. Complications: bacterial superinfection (impetigo, cellulitis), pneumonia (adults, immunocompromised), encephalitis, Reye syndrome (with aspirin), congenital varicella syndrome (pregnancy), and neonatal varicella. Diagnosis: clinical; PCR of vesicle fluid. Management: supportive; acyclovir for adolescents/adults, immunocompromised, and severe disease; varicella vaccine prevents disease; varicella-zoster immunoglobulin (VZIG) for post-exposure prophylaxis in high-risk. Reactivation causes herpes zoster (shingles).

\medterm{Variola}-- See Smallpox.

\medterm{Venereal Diseases}-- See Sexually Transmitted Infections (STIs).

\medterm{Viral Infections}-- Illnesses caused by viruses, which are intracellular pathogens that hijack host cellular machinery to replicate. They range from self-limited conditions such as the common cold and gastroenteritis to serious diseases including influenza, hepatitis, HIV, and COVID-19. Treatment is primarily supportive for most viral infections, though antiviral medications are available for specific viruses, and vaccination remains the most effective strategy for preventing many viral diseases.

\medterm{Weil's Disease}-- See Leptospirosis.

\medterm{West Nile Virus}-- A mosquito-borne flavivirus, endemic in the US, Africa, Europe, and Asia. Most infections (80\%) are asymptomatic. West Nile fever (20\%): fever, headache, myalgia, arthralgia, rash. Neuroinvasive disease (<1\%): meningitis, encephalitis, acute flaccid paralysis (anterior horn cell involvement). Risk factors for severe disease: age >50, immunosuppression, diabetes, hypertension. Diagnosis: serum/CSF WNV IgM, RT-PCR. No specific antiviral; supportive care. Prevention: mosquito avoidance and vector control.

\medterm{Whipworm (Trichuriasis)}-- Intestinal infection caused by the soil-transmitted helminth Trichuris trichiura (whipworm), endemic in tropical and subtropical regions with poor sanitation. Humans acquire infection by ingesting embryonated eggs from contaminated soil, food, or water. Larvae hatch in the small intestine and migrate to the large intestine (caecum and colon), where adult worms (3-5 cm) embed their anterior ends into the mucosa. Light infections are often asymptomatic. Heavy infections cause abdominal pain, diarrhea (sometimes bloody), tenesmus, rectal prolapse (especially in children with massive worm burden), anemia (iron-deficiency from chronic blood loss), and growth retardation. Whipworm potentiates the severity of other infections (malaria, bacterial dysentery). Diagnosis: microscopic identification of characteristic bipolar-plugged eggs in stool (Kato-Katz technique, sedimentation). Treatment: albendazole (400 mg daily for 3 days) or mebendazole (100 mg twice daily for 3 days); ivermectin is less effective alone but synergistic in combination. Mass drug administration (MDA) with albendazole targets whipworm in deworming programmes. Prevention: improved sanitation, hand washing, and avoiding soil-contaminated food.

\medterm{Yaws (Frambesia)}-- A chronic, non-venereal treponematosis caused by Treponema pallidum subspecies pertenue, closely related to the syphilis spirochete (T. pallidum pallidum). Yaws is endemic in warm, humid tropical regions of Africa, Southeast Asia, and the Pacific Islands, affecting primarily children aged 2-15 years living in rural, resource-limited areas. Transmission: direct skin-to-skin contact with an infectious lesion (abrasion, cut, or insect bite), not sexually transmitted and not vertically transmitted. Yaws is classified into stages based on clinical progression: primary yaws (mother yaw): a single, painless, papillomatous papule at the inoculation site (typically on the lower extremities, buttocks, or face) that enlarges into a raspberry-like lesion (frambesia, from French framboise = raspberry) with a characteristic yellow crust; heals spontaneously over 3-6 months, leaving a hypopigmented scar. Secondary yaws (disseminated): occurs weeks to months after the primary lesion, consisting of multiple, painless, papillomatous or ulcerated lesions (daughter yaws) on the trunk, extremities, perioral and perianal areas; hyperkeratosis of the palms and soles (crab yaws or dry yaws — painful, fissured hyperkeratotic plaques causing a crab-like gait), periostitis, and lymphadenopathy. Late/latent yaws: occurs in 10\% of untreated individuals, after 5-10 years of latency, causing destructive gummatous lesions of the skin and bones (gondou/bossed nose — hypertrophic periostitis of the maxilla), palatal and nasopharyngeal destruction (gangosa/rhinopharyngitis mutilans — destructive ulceration of the nose, palate, and pharynx), and juxta-articular nodules. Diagnosis: dark-field microscopy of serous exudate from early lesions shows motile spirochetes; serology (VDRL, RPR, TPHA, FTA-ABS — identical to syphilis, cannot distinguish between T. pallidum subspecies). Treatment: single-dose intramuscular benzathine penicillin G (1.2 million units for adults, 0.6 million units for children <10 years) is highly effective (>95\% cure rate). Alternative: oral azithromycin (single 30 mg/kg dose, maximum 2 g) — azithromycin-based mass drug administration (MDA) is the cornerstone of the WHO yaws eradication campaign. Prevention: MDA, improved hygiene, and contact tracing.

\medterm{Yellow Fever}-- A mosquito-borne viral hemorrhagic fever caused by yellow fever virus (Flavivirus), endemic in tropical Africa and South America. Incubation 3-6 days. Presents with fever, myalgia, headache, and vomiting. A proportion of patients progress to toxic phase with jaundice, hemorrhagic fever, and multi-organ failure (hepatorenal). Case fatality rate 20-60\% in severe cases. Diagnosis: RT-PCR, serology. No specific antiviral treatment; supportive care. Prevention: live-attenuated yellow fever vaccine (highly effective, required for travel to endemic areas).

\medterm{Yersinia Enterocolitica}-- A gram-negative coccobacillus of the family Enterobacteriaceae, a common cause of enterocolitis and mesenteric lymphadenitis, especially in children and young adults in temperate climates. Transmission: ingestion of contaminated food (undercooked pork — the most common source, unpasteurised milk, contaminated water), direct contact with infected animals (pigs are the primary reservoir), and person-to-person spread (rare). Pathogenicity: Y. enterocolitica carries a 70-kb virulence plasmid (pYV) encoding a type III secretion system (T3SS, injectisome) that injects Yersinia outer proteins (Yops) into host cells, disrupting actin cytoskeleton, inhibiting phagocytosis, and triggering apoptosis of macrophages. Clinical: acute enterocolitis (watery or bloody diarrhea, fever, abdominal pain, vomiting, most common in young children, self-limited over 5-14 days), pseudoappendicitis syndrome (mesenteric lymphadenitis and terminal ileitis mimicking acute appendicitis — fever, right lower quadrant pain, leukocytosis, common in older children and adolescents), and post-infectious complications (reactive arthritis, erythema nodosum, and uveitis, associated with HLA-B27). Diagnosis: stool culture (requires selective media — cefsulodin-irgasan-novobiocin/CIN agar), serology (antibodies to Yersinia antigens). Treatment: self-limited in immunocompetent patients; supportive care (rehydration) is the mainstay. Antibiotics are indicated for severe, prolonged, or invasive infections and immunocompromised patients: doxycycline + aminoglycoside, TMP-SMX, fluoroquinolones (ciprofloxacin), or third-generation cephalosporins (ceftriaxone). Y. enterocolitica produces beta-lactamase and is resistant to ampicillin, amoxicillin, and first-generation cephalosporins.

\medterm{Zika}-- Zika virus: a mosquito-borne (Aedes) flavivirus, also sexually transmitted. Usually asymptomatic or mild (fever, rash, conjunctivitis, arthralgia). Major concern: congenital Zika syndrome if acquired in pregnancy—microcephaly, brain abnormalities, eye defects; also Guillain-Barré syndrome. Diagnosis: RT-PCR (acute), IgM serology (cross-reacts with dengue); travel history. Management: supportive (no specific antiviral); prevention: mosquito avoidance, condoms, and travel advice for pregnant women; no vaccine widely available. Surveillance and vector control are central.

\medterm{Zoonoses}-- Infectious diseases that are naturally transmitted between vertebrate animals and humans. Zoonotic infections account for approximately 60\% of all human infectious diseases and 75\% of emerging infectious diseases (including HIV, SARS, MERS, Ebola, Nipah, and COVID-19). Classification by transmission route: direct zoonoses (transmission by direct contact — rabies from animal bites, anthrax from handling infected animal products), vector-borne zoonoses (transmitted by arthropod vectors — Lyme disease/Borrelia from ticks, West Nile fever from mosquitoes, plague/Yersinia pestis from fleas), food-borne zoonoses (ingestion of contaminated animal products — salmonellosis from eggs/poultry, campylobacteriosis from undercooked poultry, brucellosis from unpasteurised dairy, toxoplasmosis from undercooked meat), and water-borne zoonoses (leptospirosis from water contaminated with rodent urine). Reservoir hosts: mammals (rodents — hantavirus, plague, leptospirosis; bats — rabies, Nipah, Hendra, Ebola, Marburg, SARS, MERS; primates — HIV, yellow fever, monkeypox; pigs — influenza, Nipah, Japanese encephalitis), birds (avian influenza, West Nile virus, psittacosis), and arthropods (ticks, mosquitoes, fleas). One Health approach: integrates human health, animal health, and environmental health for surveillance, prevention, and control of zoonotic infections. Prevention: animal vaccination (rabies in dogs, brucellosis in cattle), vector control (mosquito nets, insect repellents, insecticide spraying, environmental management), food safety (pasteurisation, thorough cooking), personal protective equipment (for veterinarians, slaughterhouse workers, laboratory personnel), and surveillance of both animal and human populations for early detection of emerging zoonoses., endemic in tropical and subtropical regions with poor sanitation. Humans acquire infection by ingesting embryonated eggs from contaminated soil, food, or water. Larvae hatch in the small intestine and migrate to the large intestine (caecum and colon), where adult worms (3-5 cm) embed their anterior ends into the mucosa. Light infections are often asymptomatic. Heavy infections cause abdominal pain, diarrhea (sometimes bloody), tenesmus, rectal prolapse (especially in children with massive worm burden), anemia (iron-deficiency from chronic blood loss), and growth retardation. Whipworm potentiates the severity of other infections (malaria, bacterial dysentery). Diagnosis: microscopic identification of characteristic bipolar-plugged eggs in stool (Kato-Katz technique, sedimentation). Treatment: albendazole (400 mg daily for 3 days) or mebendazole (100 mg twice daily for 3 days); ivermectin is less effective alone but synergistic in combination. Mass drug administration (MDA) with albendazole targets whipworm in deworming programmes. Prevention: improved sanitation, hand washing, and avoiding soil-contaminated food.
